%% ============================================================================
%% Chapter 10: Idempotency, Caching & Concurrency Control
%% Mastering APIs: From Zero to Production Guru
%% ============================================================================
\chapter{Idempotency, Caching \& Concurrency Control}
\label{ch:idempotency}

%% ----------------------------------------------------------------------------
%% Executive Summary
%% ----------------------------------------------------------------------------
Networks fail in the middle. Not before your request, not after the response
arrives, but precisely in the interval where neither party can tell whether
the other side acted. Every production API lives inside that interval: a
client that times out cannot distinguish \emph{the server never received the
request} from \emph{the server executed it and the response was lost}. Three
disciplines together turn that ambiguity from a source of corruption into an
engineering property you can prove things about. \emph{Idempotency} makes
repeating an operation safe. \emph{Caching} makes repeating a read cheap.
\emph{Concurrency control} makes simultaneous writes correct. This chapter
treats all three as one coherent correctness program, because in production
they are never exercised separately: the retry that saves your client from a
timeout is the same mechanism that double-charges your customer when the
endpoint is not idempotent, and the cache that absorbs your read traffic is
the same mechanism that serves a stale balance when validation is sloppy.

We begin formally. Idempotency is defined as a fixed-point property on
observable state --- $f(f(x)) = f(x)$ --- and we map that definition onto the
HTTP method vocabulary of RFC~9110, separating what the specification
guarantees (GET, PUT, DELETE) from what well-designed APIs \emph{add}
(idempotent POST via idempotency keys)~\cite{rfc9110}. The IETF
\texttt{Idempotency-Key} draft is studied as a wire protocol with precise
replay semantics: same key and same payload replays the stored response,
same key and different payload is a 422, a concurrent duplicate in flight is
a 409, and keys expire on a 24-hour window in the style Stripe popularized.
We then dismantle the exactly-once myth: the Two Generals' Problem makes
distributed exactly-once delivery impossible, and the achievable substitute
--- at-least-once delivery combined with idempotent processing --- is what
the literature calls \emph{effectively-once}~\cite{kleppmann2017ddia}.

The caching half of the chapter is a guided tour of RFC~9111: the
\texttt{Cache-Control} directive vocabulary (\texttt{max-age}, the widely
confused \texttt{no-cache} versus \texttt{no-store}, \texttt{private} versus
\texttt{public}, \texttt{stale-while-revalidate}, \texttt{stale-if-error},
\texttt{must-revalidate}), strong and weak validators, the
\texttt{If-None-Match}/304 dance, heuristic freshness, surrogate keys for
CDN purging, \texttt{Vary}, and cache poisoning through unkeyed
inputs~\cite{rfc9110}. Concurrency control gets equal rigor: the lost-update
problem is developed as a worked interleaving, optimistic locking with
\texttt{If-Match}/ETag is given its precise status-code algebra (428, 412,
and when 409 is the honest answer), pessimistic patterns are scoped to where
they belong, and distributed locks are presented together with their failure
modes and the fencing-token fix~\cite{kleppmann2017ddia,nygard2018releaseit}.

The code is production-shaped and entirely offline. You will build a
SQLite-backed \texttt{Idempotency-Key} middleware for FastAPI and prove,
with fifty concurrent duplicate submits released on a single barrier, that
exactly one mutation occurs; a conditional-requests demo that measures the
bandwidth a 304 saves; an optimistic-locked inventory resource for the
fictional Northwind Robotics warehouse; a double-submit storm harness that
reproduces the classic double-charge incident and then cures it; and a pure
stdlib cache-policy simulator that makes RFC~9111 directive semantics
executable. Everything is typed Python~3.11, deterministic, and tested with
in-process ASGI clients~\cite{fastapidocs,pydanticdocs}.

\begin{keyconceptbox}
A retry is a second attempt at a request whose first outcome is
\emph{unknown}. Idempotency converts ``unknown'' into ``safe to repeat'':
at-least-once delivery plus idempotent processing yields effectively-once
semantics. Caching and concurrency control are the read-side and write-side
halves of the same discipline: validators (ETag, Last-Modified) let a cache
prove freshness with \texttt{If-None-Match}, and the same validators let a
writer prove it is overwriting the state it actually read with
\texttt{If-Match}. One mechanism --- the validator --- underwrites both.
\end{keyconceptbox}

%% ----------------------------------------------------------------------------
\section{Learning Objectives \& Prerequisites}
\label{sec:ch10-objectives}

After completing this chapter you will be able to:

\begin{enumerate}
  \item Define idempotency formally as a fixed-point property on observable
        state, and classify every HTTP method by its RFC~9110 safety and
        idempotency guarantees~\cite{rfc9110}.
  \item Implement the IETF \texttt{Idempotency-Key} protocol end to end:
        client key generation, fingerprinted request storage, response
        snapshotting, replay semantics, in-flight deduplication (409),
        mismatch rejection (422), expiry windows, and per-account scoping.
  \item Explain why distributed exactly-once delivery is impossible (Two
        Generals' Problem) and construct the effectively-once substitute
        from at-least-once delivery plus idempotent
        processing~\cite{kleppmann2017ddia}.
  \item Author correct \texttt{Cache-Control} policies, distinguish
        \texttt{no-cache} from \texttt{no-store}, apply
        \texttt{stale-while-revalidate} and \texttt{stale-if-error}
        deliberately, and choose strong versus weak validators correctly.
  \item Diagnose the lost-update problem from an interleaving, and prevent
        it with optimistic concurrency control using \texttt{If-Match},
        selecting 428, 412, and 409 precisely.
  \item Evaluate pessimistic locking, database constraints, and distributed
        locks (with fencing tokens) against optimistic control for a given
        workload~\cite{kleppmann2017ddia}.
  \item Prove, with a deterministic concurrency test, that a retry storm
        against an idempotent endpoint produces exactly one observable
        mutation.
\end{enumerate}

\textbf{Prerequisites.} Chapters~0--6: the HTTP wire protocol including
methods, status codes, and header semantics (Chapter~1); consuming APIs from
Python with \texttt{httpx} and retry discipline (Chapter~2); REST resource
modeling (Chapter~3); JSON serialization and canonicalization
(Chapter~4); the FastAPI application factory, dependency injection,
middleware, and \texttt{TestClient} patterns (Chapter~6). Comfort with
\texttt{threading.Barrier}, \texttt{sqlite3}, and \texttt{hashlib} from the
Python standard library is assumed; no database server, network access, or
cloud account is required anywhere in this chapter.

\begin{table}[h]
  \centering
  \small
  \begin{tabularx}{\textwidth}{l X}
    \toprule
    \textbf{Prerequisite} & \textbf{Where it is used here} \\
    \midrule
    HTTP methods, status codes, headers & method idempotency algebra; 304/409/412/422/428 selection \\
    Retry loops with backoff (Ch.~2) & why non-idempotent retry is unsafe; the money-transfer argument \\
    REST resources and representations (Ch.~3) & ETags as representation validators; conditional PUT \\
    JSON canonicalization (Ch.~4) & request fingerprinting for idempotency keys \\
    FastAPI middleware and TestClient (Ch.~6) & the idempotency middleware and all offline tests \\
    \bottomrule
  \end{tabularx}
  \caption{Prerequisite map for Chapter~10.}
  \label{tab:ch10-prereqs}
\end{table}

%% ----------------------------------------------------------------------------
\section{Deep Technical Mechanics \& Architectural Concepts}
\label{sec:ch10-mechanics}

\subsection{Idempotency: A Formal Definition}
\label{sec:ch10-idempotency-formal}

\begin{definition}[Idempotent operation]
\label{def:idempotent}
Let $S$ be the set of observable server states and let an API operation be a
function $f : S \to S$. The operation is \emph{idempotent} iff applying it
twice is observationally indistinguishable from applying it once:
\[
  f(f(x)) = f(x) \quad \text{for all } x \in S.
\]
``Observable'' deliberately includes side effects a client can detect ---
ledger entries, emails sent, robots dispatched --- and excludes
implementation-internal artifacts such as log lines or monotonic request
counters.
\end{definition}

Two consequences deserve emphasis. First, idempotency is a property of the
\emph{effect}, not of the \emph{response}: a replayed request may legally
return the freshly stored receipt rather than re-executing, and it may even
return a different status code in degenerate cases (a second DELETE
returning 404) while remaining idempotent, because the \emph{state} is
unchanged. Second, idempotency is not safety. A safe method (GET, HEAD)
changes nothing at all; an idempotent method may change state on the first
application but not on repetition. The identity function is the only safe
operation; every fixed-point operation is idempotent.

RFC~9110 assigns each standard method a fixed position in this
algebra~\cite{rfc9110}, reproduced in Table~\ref{tab:ch10-methods}. The
subtle row is DELETE: the first call transitions the resource to
``absent''; every subsequent call finds the fixed point already reached.
POST is the only mainstream method with \emph{no} guarantee --- which is
precisely why POST is the method that carries charges, bookings, and
dispatches, and precisely why the industry layered idempotency keys on top
of it.

\begin{table}[h]
  \centering
  \small
  \begin{tabularx}{\textwidth}{l c c X}
    \toprule
    \textbf{Method} & \textbf{Safe} & \textbf{Idempotent} & \textbf{Effect of a retry} \\
    \midrule
    GET, HEAD & yes & yes & none; caches may even answer locally \\
    OPTIONS, TRACE & yes & yes & none \\
    PUT & no & yes & same absolute state written again \\
    DELETE & no & yes & resource remains absent (fixed point) \\
    POST & no & \textbf{no} & a second, independent side effect \\
    PATCH & no & no (in general) & depends on the delta encoding \\
    \bottomrule
  \end{tabularx}
  \caption{Safety and idempotency per RFC~9110~\cite{rfc9110}. PATCH is
  idempotent only if the patch document encodes an absolute, not relative,
  transformation.}
  \label{tab:ch10-methods}
\end{table}

\begin{examplebox}
\textbf{Fixed points in one inventory endpoint.} Consider
\texttt{PUT /inventory/NWR-AXLE-42} with body
\texttt{\{"quantity": 30\}}. Whether executed once or fifty times, the
observable state ends at \texttt{quantity = 30}: the operation is
idempotent. Now consider \texttt{POST /inventory/NWR-AXLE-42/adjustments}
with body \texttt{\{"delta": -5\}}. Each execution appends a new adjustment;
fifty executions reserve 250 units. The POST is not idempotent, and no
amount of careful client retry logic can make it safe unless the server
adds a deduplication mechanism. That mechanism is the idempotency key.
\end{examplebox}

\subsection{The Idempotency-Key Protocol}
\label{sec:ch10-idempotency-keys}

The \texttt{Idempotency-Key} pattern, standardized as an IETF draft and
popularized by Stripe, upgrades an unsafe POST into an idempotent operation
by attaching a client-generated unique key to each logical request. The
server stores the outcome keyed by that value and answers duplicates from
the store instead of re-executing. The protocol has five design decisions,
each of which is a failure mode when gotten wrong.

\textbf{Key generation.} The key must be generated by the \emph{client}, one
per logical operation (typically a UUIDv4 minted when the user opens the
checkout form, not when the request is sent). All retries of that operation
--- manual retries, library retries, load-balancer retries --- reuse the
same key. A new key means a new operation; that is the entire contract.

\textbf{Storage design.} The server persists, per key: a \emph{request
fingerprint} (a SHA-256 over method, route, and canonical request body), a
\emph{response snapshot} (status code, body, content type), a lifecycle
state (\texttt{in\_flight} or \texttt{completed}), and a creation timestamp
for expiry. The fingerprint is what makes key reuse detectable: it must
cover the route, not only the body, or the same JSON sent to two endpoints
will collide.

\textbf{Replay semantics.} Three collision outcomes, three status codes:
\begin{itemize}
  \item Same key, same fingerprint, state \texttt{completed} $\rightarrow$
        replay the stored response byte-for-byte, marked with a header such
        as \texttt{Idempotency-Replayed: true} so observability can
        distinguish replays from executions.
  \item Same key, different fingerprint $\rightarrow$ \textbf{422}. The
        client reused a key for a different logical operation; replaying
        would silently corrupt intent, and executing would double-apply.
  \item Same key, same fingerprint, state \texttt{in\_flight}
        $\rightarrow$ \textbf{409} with \texttt{Retry-After}, because the
        original is still executing and its outcome does not exist yet.
        (Some drafts discuss 425 Too Early; 409 is the widely deployed
        choice because the condition is a genuine state conflict.)
\end{itemize}

\textbf{Expiry.} Keys cannot be stored forever; Stripe retains them for 24
hours. The expiry window is a correctness parameter, not housekeeping: it
bounds how long a duplicate can be recognized as a duplicate, and therefore
bounds how late a retry may arrive. Choose it from your client's worst-case
retry horizon, not from storage cost.

\textbf{Scoping.} Keys are unique per \emph{account}, not globally --- two
tenants must never replay each other's responses, and key uniqueness
pressure should not cross tenant boundaries. Endpoints are bound implicitly
by the fingerprint (a key reused on a different route mismatches), which
prevents cross-endpoint replay without extra bookkeeping.

Figure~\ref{fig:ch10-replay-flow} shows the full decision flow as
implemented in Listing~10.1.

\begin{figure}[h]
\centering
\begin{tikzpicture}[
    font=\small,
    startstop/.style={rectangle, rounded corners, draw=packtgray, fill=packtorange!15,
                      text width=3.3cm, align=center, minimum height=0.9cm},
    decision/.style={diamond, draw=packtgray, fill=blue!8, text width=2.6cm,
                     align=center, inner sep=1pt, aspect=2.2},
    outcome/.style={rectangle, draw=packtgray, fill=green!10, text width=3.1cm,
                    align=center, minimum height=0.9cm},
    flow/.style={->, thick, packtgray},
    node distance=0.55cm and 0.9cm]

  \node[startstop] (arrive) {POST arrives with \texttt{Idempotency-Key}};
  \node[decision, below=of arrive] (seen) {Key seen before?};
  \node[decision, below left=0.7cm and -0.2cm of seen] (match) {Fingerprint matches?};
  \node[decision, below right=0.7cm and -0.6cm of seen] (state) {State \texttt{completed}?};
  \node[outcome, below=of match] (unproc) {\textbf{422} key reused with different payload};
  \node[outcome, below=of state] (inflight) {\textbf{409} + \texttt{Retry-After}: still in flight};
  \node[outcome, right=1.0cm of inflight] (replay) {Replay stored response + \texttt{Idempotency-Replayed}};
  \node[startstop, right=1.0cm of replay] (execute) {Mark \texttt{in\_flight}, execute, snapshot response};

  \draw[flow] (arrive) -- (seen);
  \draw[flow] (seen) -| node[note, pos=0.25, above]{yes, and fingerprint differs} (match);
  \draw[flow] (match) -- node[note, left]{no} (unproc);
  \draw[flow] (seen) -| node[note, pos=0.3, above]{yes, same payload} (state);
  \draw[flow] (state) -- node[note, left]{no} (inflight);
  \draw[flow] (state) -- node[note, above]{yes} (replay);
  \draw[flow] (seen.east) -| node[note, pos=0.2, above right]{no (or expired)} (execute.north);
\end{tikzpicture}
\caption{Idempotency-key decision flow. Exactly one path executes business
logic; every other path returns a deterministic, side-effect-free answer.}
\label{fig:ch10-replay-flow}
\end{figure}

\begin{warningbox}
The response snapshot must be written in the same transaction as the state
mutation, or the system has a crash window in which the effect happened but
the receipt was never recorded. A retry arriving in that window either
double-executes (if the key row was rolled back) or deadlocks (if the key
row survived as \texttt{in\_flight} forever). Store first, commit once,
replay forever --- and never store the response before the mutation is
durable.
\end{warningbox}

\subsection{The Exactly-Once Illusion}
\label{sec:ch10-exactly-once}

Every distributed systems curriculum eventually confronts the Two Generals'
Problem: two armies must agree on an attack time using messengers who may be
captured. No finite protocol of unreliable messages can guarantee both sides
act together, because the last messenger's arrival is itself unknowable to
the sender. An HTTP request over a flaky network \emph{is} the Two Generals:
after a timeout, the client cannot know whether the request was lost before
execution or the response was lost after it. Exactly-once delivery is
therefore not a difficult engineering goal; it is a theorem-level
impossibility~\cite{kleppmann2017ddia}.

What is achievable --- and what production systems actually ship --- is the
composition
\[
  \text{at-least-once delivery} \;+\; \text{idempotent processing}
  \;=\; \text{effectively-once semantics}.
\]
The transport is allowed to deliver duplicates (and will); the processor is
built so duplicates are observationally inert. Every deduplication mechanism
has a \emph{window}: the idempotency-key TTL, the message-ID retention
period, the dedup log's compaction horizon. A duplicate arriving outside the
window is indistinguishable from a new request. Sizing the window from the
client's maximum retry horizon, and monitoring how full it runs, is the
operational discipline that makes ``effectively'' honest.

\subsection{HTTP Caching Machinery per RFC 9110/9111}
\label{sec:ch10-caching}

HTTP caching exists to eliminate two things: \emph{latency} (a cache closer
than the origin answers faster) and \emph{origin load} (a shared cache
answers instead of the origin at all). Both are achieved by the same small
machinery: a \emph{freshness lifetime} the origin assigns, and
\emph{validators} the cache uses to prove freshness when the lifetime has
run out~\cite{rfc9110}.

\textbf{Freshness directives.} \texttt{Cache-Control} is a comma-separated
directive vocabulary; the directives that matter most are summarized in
Table~\ref{tab:ch10-directives}. The single most expensive confusion in the
vocabulary is \texttt{no-cache} versus \texttt{no-store}. Despite its name,
\texttt{no-cache} permits \emph{storage} --- it forbids \emph{serving}
without revalidation. \texttt{no-store} forbids storage entirely and is the
correct choice for responses containing credentials, account data, or
one-time tokens.

\begin{table}[h]
  \centering
  \small
  \begin{tabularx}{\textwidth}{l X}
    \toprule
    \textbf{Directive} & \textbf{Semantics} \\
    \midrule
    \texttt{max-age=n} & fresh for $n$ seconds; a cache may serve without contacting the origin \\
    \texttt{no-cache} & may store, \emph{must} revalidate before every reuse \\
    \texttt{no-store} & may not store at all; every request reaches the origin \\
    \texttt{private} & only the single-user (browser) cache may store \\
    \texttt{public} & shared caches (CDNs, proxies) may store \\
    \texttt{stale-while-revalidate=n} & may serve stale for $n$ seconds past expiry while refreshing in the background \\
    \texttt{stale-if-error=n} & may serve stale for $n$ seconds if the origin errors \\
    \texttt{must-revalidate} & once stale, never serve without a successful revalidation (overrides stale serving) \\
    \bottomrule
  \end{tabularx}
  \caption{Core \texttt{Cache-Control} response directives (RFC~9111).}
  \label{tab:ch10-directives}
\end{table}

\textbf{Validators: strong and weak.} When a stored response is stale, the
cache revalidates with a \emph{conditional request}. Two validators exist.
\texttt{Last-Modified} is a one-second-granularity timestamp ---
\emph{weak}, because two changes inside one second are invisible to it.
\texttt{ETag} is an opaque tag; a \emph{strong} ETag must change whenever
the representation changes by a single byte, while a weak ETag (prefixed
\texttt{W/}) asserts only semantic equivalence. Strong ETags support range
requests and byte-sensitive caches; weak ETags are for representations whose
bytes may legitimately vary (e.g.\ recompressed but semantically identical
bodies). The golden rule: a strong ETag is a content address --- compute it
from the representation bytes (SHA-256 is ideal) or from a version counter
when the representation is a pure function of the versioned state.
\emph{Never} let a client invent the validator.

The conditional flow is two headers deep: the cache presents its stored
validator with \texttt{If-None-Match} (or \texttt{If-Modified-Since}), and
the origin answers \textbf{304 Not Modified} with no body when the
validator still matches. Figure~\ref{fig:ch10-conditional} shows the
sequence; Listing~10.2 measures the saved bytes.

\begin{figure}[h]
\centering
\begin{tikzpicture}[
    font=\small,
    actor/.style={rectangle, draw=packtgray, fill=blue!8, minimum width=2.2cm,
                  minimum height=0.8cm},
    msg/.style={->, thick, packtgray},
    rmsg/.style={->, thick, dashed, packtgray},
    node distance=2.6cm]

  \node[actor] (client) {Client};
  \node[actor, right=of client] (cache) {Cache / CDN};
  \node[actor, right=of cache] (origin) {Origin API};
  \draw[packtgray] ($(client.south west)+(0.2,-0.1)$) -- ++(0,-5.6);
  \draw[packtgray] ($(cache.south west)+(0.2,-0.1)$) -- ++(0,-5.6);
  \draw[packtgray] ($(origin.south west)+(0.2,-0.1)$) -- ++(0,-5.6);

  \draw[msg]  ($(client.south)+(0.2,-0.35)$) -- node[above, note]{GET /parts/NWR-AXLE-42} ($(cache.south)+(-2.2,-0.35)$);
  \draw[msg]  ($(cache.south)+(0.2,-0.9)$)  -- node[above, note]{GET (cache miss)} ($(origin.south)+(-2.2,-0.9)$);
  \draw[rmsg] ($(origin.south)+(-2.2,-1.5)$) -- node[above, note]{200 + body + \texttt{ETag: "a1\ldots"} + \texttt{max-age=60}} ($(cache.south)+(0.2,-1.5)$);
  \draw[rmsg] ($(cache.south)+(-2.2,-2.05)$) -- node[above, note]{200 + body (stored)} ($(client.south)+(0.2,-2.05)$);

  \draw[msg]  ($(client.south)+(0.2,-3.0)$) -- node[above, note]{GET (after expiry)} ($(cache.south)+(-2.2,-3.0)$);
  \draw[msg]  ($(cache.south)+(0.2,-3.55)$) -- node[above, note]{GET + \texttt{If-None-Match: "a1\ldots"}} ($(origin.south)+(-2.2,-3.55)$);
  \draw[rmsg] ($(origin.south)+(-2.2,-4.15)$) -- node[above, note]{\textbf{304 Not Modified} (no body, freshness reset)} ($(cache.south)+(0.2,-4.15)$);
  \draw[rmsg] ($(cache.south)+(-2.2,-4.7)$) -- node[above, note]{200 from stored body} ($(client.south)+(0.2,-4.7)$);
\end{tikzpicture}
\caption{Conditional request sequence. The revalidation round trip still
costs latency, but the 304 transfers zero representation bytes; with
\texttt{max-age} still fresh, the round trip disappears entirely.}
\label{fig:ch10-conditional}
\end{figure}

\textbf{Heuristic freshness.} When the origin sends validators but no
explicit lifetime, caches may apply the RFC~9111 heuristic:
$\text{freshness} = 0.1 \times (\text{Date} - \text{Last\text{-}Modified})$.
A document untouched for a week is guessed fresh for about 17 hours. The
heuristic is a fallback, not a design; explicit \texttt{max-age} is always
preferable because it is a decision rather than a guess.

\textbf{Surrogate keys and CDNs.} A surrogate key (e.g.\
\texttt{Surrogate-Key: part:NWR-AXLE-42 catalog}) tags stored responses at
the CDN so that one mutation can purge thousands of URLs by tag instead of
enumerating paths. Tag-based invalidation converts the classic ``two hard
things'' joke into an operational primitive: writes emit tags, purges
reference tags, and the cache invalidation problem becomes a bookkeeping
problem.

\textbf{Vary and cache poisoning.} The cache key is, by default, method +
URI. If the response depends on any request header --- \texttt{Accept},
\texttt{Accept-Encoding}, \texttt{Accept-Language}, an authorization scope
--- the response must declare it with \texttt{Vary}, or the cache will serve
one client's representation to another. The hostile version of this mistake
is \emph{cache poisoning}: if an unkeyed input (a reflected header, a query
parameter ignored by the key) influences the response, an attacker can store
a poisoned representation under an innocent key and have the cache deliver
it to everyone~\cite{nygard2018releaseit}. Rule: every byte of the response
must be a function of the cache key plus the \texttt{Vary} set.

\subsection{Concurrency Control: The Lost-Update Problem}
\label{sec:ch10-concurrency}

Caching controls how \emph{readers} see state; concurrency control decides
what happens when \emph{writers} collide. The canonical failure is the lost
update, worked out interleaving by interleaving in
Figure~\ref{fig:ch10-lost-update}: two clients read the same inventory row,
each computes a new value from its private copy, and the second write
silently overwrites the first. Neither client errored. The invariant
``reserved $\leq$ on hand'' is violated anyway. The bug is not in either
request; it is in the \emph{composition} --- which is why no amount of
single-request validation finds it~\cite{kleppmann2017ddia}.

\begin{figure}[h]
\centering
\begin{tikzpicture}[
    font=\small,
    step/.style={rectangle, draw=packtgray, fill=blue!8, text width=4.4cm,
                 align=left, minimum height=0.85cm},
    badstep/.style={rectangle, draw=packtgray, fill=packtorange!20, text width=4.4cm,
                    align=left, minimum height=0.85cm},
    timeline/.style={->, thick, packtgray},
    node distance=0.5cm]

  \node[flowstep] (a1) {A: GET $\rightarrow$ \texttt{qty=35}, \texttt{ETag "v1"}};
  \node[flowstep, below=of a1] (b1) {B: GET $\rightarrow$ \texttt{qty=35}, \texttt{ETag "v1"}};
  \node[flowstep, below=of b1] (a2) {A: reserves 5 $\rightarrow$ PUT \texttt{qty=30}, \texttt{If-Match "v1"} \checkmark};
  \node[badstep, below=of a2] (b2) {B: reserves 8 $\rightarrow$ PUT \texttt{qty=27}, \texttt{If-Match "v1"}};

  \node[left=0.6cm of a1, font=\footnotesize\itshape, text=packtgray] {time};
  \draw[timeline] ($(a1.north west)+(-0.35,0.25)$) -- ($(b2.south west)+(-0.35,-0.25)$);

  \node[badstep, right=1.1cm of b2, text width=5.3cm] (verdict)
    {\textbf{Without} \texttt{If-Match}: B's write lands, A's reservation of
     5 is lost; 43 units promised against 35 on hand. \textbf{With}
     \texttt{If-Match}: B's precondition fails $\rightarrow$ 412, B re-reads
     \texttt{qty=30}, reserves 8 $\rightarrow$ \texttt{qty=22}. Nothing
     lost.};
  \draw[->, thick, packtgray, dashed] (b2) -- (verdict);
\end{tikzpicture}
\caption{The lost-update interleaving. Reads interleave harmlessly; the
second blind write is the corruption. A conditional write converts the race
into a deterministic rejection.}
\label{fig:ch10-lost-update}
\end{figure}

\textbf{Optimistic locking with If-Match.} HTTP's native answer is the
conditional write. The server stamps every representation with a strong ETag
derived from a version counter; a writer echoes it in \texttt{If-Match}; the
server applies the write inside an atomic compare-and-set:
\[
  \text{apply} \iff \text{version}_{\text{current}} = \text{version}_{\text{sent}}.
\]
The status-code algebra is precise and worth memorizing:
\begin{itemize}
  \item \textbf{428 Precondition Required} --- the server \emph{requires}
        conditional writes and the client sent none. This is the API
        declaring ``blind overwrites are forbidden here.''
  \item \textbf{412 Precondition Failed} --- the client sent a precondition
        and it evaluated false: someone wrote first. The correct client
        behavior is read--re-decide--retry.
  \item \textbf{409 Conflict} --- the request conflicts with server state
        for \emph{business} reasons (duplicate natural key, invalid state
        transition), not with a representation version. Use 409 for semantic
        conflicts and in-flight idempotency collisions; use 412 strictly for
        failed preconditions. Confusing them sends clients down the wrong
        recovery path.
\end{itemize}
Optimistic control shines when contention is rare: zero locks held, no
blocking, and the retry cost is paid only by actual collisions. Its price is
the read-modify-write retry loop, which under hot contention can livelock;
the fix is bounded retries with jittered backoff and escalation.

\textbf{Pessimistic patterns.} When contention is hot or the retry of a
failed operation is expensive (a payment re-authorization, a compiled
report), take the lock up front: \texttt{SELECT \ldots\ FOR UPDATE} in the
database, or a serialized queue per aggregate. Pessimism trades throughput
for predictability; it is the right tool for short critical sections over
hot rows, and the wrong tool for long user think-times, where a held lock
outlives the session.

\textbf{Distributed locks and fencing.} A lock service (Redis,
ZooKeeper-style consensus) extends mutual exclusion across processes, but a
paused lock holder is indistinguishable from a dead one: the lease expires,
a second client acquires the lock, and the paused client wakes up and writes
anyway. The fix is a \emph{fencing token} --- a monotonic number issued with
every lock acquisition --- which the resource checks on every write,
rejecting any token older than the newest it has seen. The lock provides
mutual exclusion for efficiency; the fencing token provides correctness for
safety. A distributed lock without fencing is a race condition with extra
infrastructure~\cite{kleppmann2017ddia}.

\begin{table}[h]
  \centering
  \small
  \begin{tabularx}{\textwidth}{l X X}
    \toprule
    \textbf{Strategy} & \textbf{Choose when} & \textbf{Cost} \\
    \midrule
    Optimistic (If-Match/CAS) & low-to-moderate contention; cheap retries & retry loop; livelock risk on hot keys \\
    Pessimistic (DB row lock) & hot rows; expensive retries; short critical sections & blocking, lock management, deadlock risk \\
    Serialized queue / single writer & per-aggregate ordering is natural & throughput ceiling of one writer \\
    Distributed lock + fencing token & cross-process mutual exclusion & lock service dependency; fencing enforced at the resource \\
    \bottomrule
  \end{tabularx}
  \caption{Concurrency-control strategy selection.}
  \label{tab:ch10-strategies}
\end{table}

\subsection{Retries Meet Idempotency: The Money-Transfer Argument}
\label{sec:ch10-retries}

Chapter~2 taught clients to retry with exponential backoff and jitter.
Chapter~10 supplies the correctness condition that makes that teaching safe:
\emph{a client may retry freely iff the operation is idempotent or carries
an idempotency key}. Consider the canonical non-idempotent operation --- a
money transfer POST --- under a timeout:

\begin{enumerate}
  \item The client sends \texttt{POST /transfers \{amount\_cents: 2500\}}.
  \item The server executes it; the response is lost in the network.
  \item The client retries. Without a key, the server executes \emph{again}.
        The customer is charged twice; the ledger is authoritative; the
        duplicate must now be discovered and refunded by a human.
\end{enumerate}

The argument generalizes into a design rule with three clauses. (1)~The
outcome of a timed-out request is \emph{unknown}, never ``failed'' ---
treating it as failed is the bug. (2)~Retrying an unknown-outcome
non-idempotent request creates exactly one duplicate per ambiguity window;
load balancers, proxies, and HTTP libraries all create such windows, often
without the application code's knowledge. (3)~The only client-side fix is to
make the retry \emph{recognizable}: same key, same payload, so the server's
dedup machinery can collapse the storm to one execution. Idempotency keys
are therefore not a server feature; they are the second half of the retry
contract, and an SDK that retries POST without minting a key is unsafe by
construction.

%% ----------------------------------------------------------------------------
\section{Production-Ready Code Implementation}
\label{sec:ch10-code}

All listings run offline on Python~3.11+ against in-process ASGI test
clients; concurrency is exercised with \texttt{threading.Barrier} and
assertions on final state, never on wall-clock timing. Set up once per
machine from PowerShell:

\begin{minted}{text}
python -m venv .venv
.\.venv\Scripts\Activate.ps1
pip install fastapi httpx pydantic pytest
pytest -q
\end{minted}

\subsection{Listing 10.1 --- SQLite-Backed Idempotency-Key Middleware}
\label{sec:ch10-listing1}

The middleware implements the full decision flow of
Figure~\ref{fig:ch10-replay-flow}: atomic check-and-insert against a SQLite
store (the \texttt{PRIMARY KEY} constraint \emph{is} the concurrency
control), fingerprinting over method + route + canonical JSON body, response
snapshotting, 422 on fingerprint mismatch, 409 with \texttt{Retry-After}
while the twin is in flight, per-account scoping, and a Stripe-style 24-hour
TTL. Two deliberate design choices deserve attention. First, keys are
claimed \emph{before} execution and the row is \emph{abandoned} on
exceptions and 5xx responses, because a server-side failure is an outcome a
client must be allowed to retry. Second, the store's single connection is
guarded by one \texttt{RLock}; correctness lives in the atomic
\texttt{INSERT}, not in the volume of the bookkeeping.

\begin{minted}{python}
"""Listing 10.1 -- SQLite-backed Idempotency-Key middleware for FastAPI.

Northwind Robotics Payments API. Offline-first, typed Python 3.11+.
Implements the IETF `Idempotency-Key` draft semantics:

* first request with a key executes and snapshots the response,
* same key + same payload replays the stored response,
* same key + different payload is rejected with 422,
* a concurrent duplicate while the first is in flight gets 409,
* keys expire after a configurable TTL (default: 24 hours, Stripe-style),
* keys are scoped per account (header ``X-Account-Id``).
"""

from __future__ import annotations

import hashlib
import json
import sqlite3
import threading
import time
from dataclasses import dataclass
from typing import Optional

from fastapi import FastAPI, Request
from fastapi.responses import JSONResponse, Response
from pydantic import BaseModel, Field
from starlette.middleware.base import (
    BaseHTTPMiddleware,
    RequestResponseEndpoint,
)

IDEMPOTENCY_KEY_HEADER = "idempotency-key"
REPLAYED_HEADER = "idempotency-replayed"
DEFAULT_TTL_SECONDS = 24 * 60 * 60  # Stripe-style 24-hour key expiry
IDEMPOTENT_METHODS = frozenset({"POST", "PUT", "PATCH"})


# --------------------------------------------------------------------------
# Persistence: the idempotency key store
# --------------------------------------------------------------------------
@dataclass(frozen=True)
class StoredResponse:
    """Immutable snapshot of a completed response."""

    status_code: int
    body: bytes
    content_type: str


class IdempotencyStore:
    """SQLite-backed key store.

    One row per (scope, key). Rows transition ``in_flight`` -> ``completed``.
    A single connection guarded by an RLock keeps concurrent handler
    threads serialized, so ``begin`` is an atomic check-and-insert.
    """

    def __init__(self, path: str = ":memory:", ttl_seconds: float = DEFAULT_TTL_SECONDS) -> None:
        self.ttl_seconds = ttl_seconds
        self._lock = threading.RLock()
        self._conn = sqlite3.connect(path, check_same_thread=False)
        self._conn.execute(
            """
            CREATE TABLE IF NOT EXISTS idempotency_keys (
                scope        TEXT NOT NULL,
                idem_key     TEXT NOT NULL,
                fingerprint  TEXT NOT NULL,
                state        TEXT NOT NULL CHECK (state IN ('in_flight', 'completed')),
                status_code  INTEGER,
                response_body BLOB,
                content_type TEXT,
                created_at   REAL NOT NULL,
                PRIMARY KEY (scope, idem_key)
            )
            """
        )
        self._conn.commit()

    def _prune_expired(self, now: float) -> None:
        self._conn.execute(
            "DELETE FROM idempotency_keys WHERE created_at + ? < ?",
            (self.ttl_seconds, now),
        )

    def begin(
        self, scope: str, key: str, fingerprint: str, now: Optional[float] = None
    ) -> tuple[str, Optional[StoredResponse]]:
        """Atomically claim a key or classify the collision.

        Returns one of:
          ("started",  None)   -- caller won the race and must execute,
          ("replay",   stored) -- same fingerprint, completed: replay it,
          ("mismatch", None)   -- same key, different payload: 422,
          ("in_flight", None)  -- same key in flight concurrently: 409.
        """
        now = time.time() if now is None else now
        with self._lock:
            self._prune_expired(now)
            try:
                self._conn.execute(
                    "INSERT INTO idempotency_keys"
                    " (scope, idem_key, fingerprint, state, created_at)"
                    " VALUES (?, ?, ?, 'in_flight', ?)",
                    (scope, key, fingerprint, now),
                )
                self._conn.commit()
                return ("started", None)
            except sqlite3.IntegrityError:
                row = self._conn.execute(
                    "SELECT fingerprint, state, status_code, response_body,"
                    "       content_type"
                    " FROM idempotency_keys WHERE scope = ? AND idem_key = ?",
                    (scope, key),
                ).fetchone()
            if row is None:  # pragma: no cover - defensive
                raise RuntimeError("idempotency row vanished mid-transaction")
            if row[0] != fingerprint:
                return ("mismatch", None)
            if row[1] == "in_flight":
                return ("in_flight", None)
            return ("replay", StoredResponse(row[2], bytes(row[3]), row[4]))

    def complete(
        self, scope: str, key: str, stored: StoredResponse
    ) -> None:
        """Attach the finished response snapshot to the claimed key."""
        with self._lock:
            self._conn.execute(
                "UPDATE idempotency_keys"
                " SET state = 'completed', status_code = ?,"
                "     response_body = ?, content_type = ?"
                " WHERE scope = ? AND idem_key = ? AND state = 'in_flight'",
                (stored.status_code, stored.body, stored.content_type, scope, key),
            )
            self._conn.commit()

    def abandon(self, scope: str, key: str) -> None:
        """Release the key after a 5xx/exception so the client may retry."""
        with self._lock:
            self._conn.execute(
                "DELETE FROM idempotency_keys"
                " WHERE scope = ? AND idem_key = ? AND state = 'in_flight'",
                (scope, key),
            )
            self._conn.commit()

    def close(self) -> None:
        with self._lock:
            self._conn.close()


def fingerprint_request(method: str, path: str, body: bytes) -> str:
    """SHA-256 over method + route + canonical body.

    Hashing only the body is a classic bug: the same JSON payload sent to
    two different endpoints must not collide, so the route is part of the
    fingerprint. JSON bodies are canonicalized so key order is irrelevant.
    """
    try:
        canonical = json.dumps(
            json.loads(body.decode("utf-8")), sort_keys=True, separators=(",", ":")
        ).encode("utf-8")
    except (ValueError, UnicodeDecodeError):
        canonical = body
    digest = hashlib.sha256()
    digest.update(method.upper().encode("ascii"))
    digest.update(b"\x00")
    digest.update(path.encode("utf-8"))
    digest.update(b"\x00")
    digest.update(canonical)
    return digest.hexdigest()


def _problem(status: int, title: str, detail: str, type_slug: str) -> JSONResponse:
    return JSONResponse(
        status_code=status,
        media_type="application/problem+json",
        content={
            "type": f"https://api.northwind-robotics.example/problems/{type_slug}",
            "title": title,
            "status": status,
            "detail": detail,
        },
    )


class IdempotencyMiddleware(BaseHTTPMiddleware):
    """Applies Idempotency-Key semantics to unsafe methods."""

    def __init__(
        self,
        app: FastAPI,
        store: IdempotencyStore,
        methods: frozenset[str] = IDEMPOTENT_METHODS,
    ) -> None:
        super().__init__(app)
        self.store = store
        self.methods = methods

    async def dispatch(
        self, request: Request, call_next: RequestResponseEndpoint
    ) -> Response:
        key = request.headers.get(IDEMPOTENCY_KEY_HEADER)
        if request.method not in self.methods or not key:
            return await call_next(request)

        scope = request.headers.get("x-account-id", "anonymous")
        body = await request.body()
        fingerprint = fingerprint_request(request.method, request.url.path, body)

        outcome, stored = self.store.begin(scope, key, fingerprint)
        if outcome == "mismatch":
            return _problem(
                422,
                "Idempotency-key payload mismatch",
                "This Idempotency-Key was already used with a different request"
                " payload. Keys identify one logical request; mint a new key"
                " for a new operation.",
                "idempotency-key-mismatch",
            )
        if outcome == "in_flight":
            response = _problem(
                409,
                "Request with this idempotency key is in flight",
                "A concurrent request carrying the same Idempotency-Key is"
                " still being processed. Retry after a short backoff.",
                "idempotency-key-in-flight",
            )
            response.headers["Retry-After"] = "1"
            return response
        if outcome == "replay" and stored is not None:
            return Response(
                content=stored.body,
                status_code=stored.status_code,
                media_type=stored.content_type,
                headers={REPLAYED_HEADER: "true"},
            )

        # We claimed the key: execute and snapshot, or abandon on failure.
        try:
            downstream = await call_next(request)
            body_bytes = b""
            async for chunk in downstream.body_iterator:
                body_bytes += chunk
        except Exception:
            self.store.abandon(scope, key)
            raise

        content_type = downstream.headers.get("content-type", "application/json")
        if downstream.status_code < 500:
            self.store.complete(
                scope, key, StoredResponse(downstream.status_code, body_bytes, content_type)
            )
        else:
            # Server-side failures are never replayed; free the key.
            self.store.abandon(scope, key)
        return Response(
            content=body_bytes,
            status_code=downstream.status_code,
            media_type=content_type,
        )


# --------------------------------------------------------------------------
# Demo application: Northwind Robotics payments ledger
# --------------------------------------------------------------------------
class TransferRequest(BaseModel):
    """One funds movement between internal accounts."""

    from_account: str = Field(min_length=1)
    to_account: str = Field(min_length=1)
    amount_cents: int = Field(gt=0)


class TransferReceipt(BaseModel):
    transfer_id: str
    from_account: str
    to_account: str
    amount_cents: int
    sequence: int


class Ledger:
    """Thread-safe in-memory ledger; the observable side effect under test."""

    def __init__(self) -> None:
        self._lock = threading.Lock()
        self._transfers: list[TransferReceipt] = []

    def post(self, request: TransferRequest) -> TransferReceipt:
        with self._lock:
            receipt = TransferReceipt(
                transfer_id=f"tr_{len(self._transfers) + 1:06d}",
                from_account=request.from_account,
                to_account=request.to_account,
                amount_cents=request.amount_cents,
                sequence=len(self._transfers) + 1,
            )
            self._transfers.append(receipt)
            return receipt

    def count(self) -> int:
        with self._lock:
            return len(self._transfers)

    def total_cents(self) -> int:
        with self._lock:
            return sum(t.amount_cents for t in self._transfers)


def create_app(store: Optional[IdempotencyStore] = None) -> FastAPI:
    """Application factory; each test gets an isolated store + ledger."""
    app = FastAPI(title="Northwind Robotics Payments", version="1.0.0")
    app.state.store = store if store is not None else IdempotencyStore()
    app.state.ledger = Ledger()
    app.add_middleware(IdempotencyMiddleware, store=app.state.store)

    @app.post("/api/v1/payments/transfers", status_code=201)
    def create_transfer(payload: TransferRequest) -> TransferReceipt:
        return app.state.ledger.post(payload)

    return app
\end{minted}

The test suite proves every branch of the decision flow, culminating in the
50-thread duplicate storm: fifty workers released on one barrier submit the
identical keyed payload; the assertion battery checks that the ledger holds
exactly one transfer, that all fifty callers received the byte-identical
receipt, and that exactly one response was a non-replayed origin execution.

\begin{minted}{python}
"""Listing 10.1 test suite -- replay, mismatch, scope, and 50-thread storm.

Concurrency tests assert on final state and response equality only; a
threading.Barrier maximizes overlap without any timing dependence.
"""

from __future__ import annotations

import threading

import pytest
from fastapi.testclient import TestClient

from listing_10_1_idempotency_middleware import (
    REPLAYED_HEADER,
    IdempotencyStore,
    create_app,
    fingerprint_request,
)

PAYLOAD = {
    "from_account": "ops-checking",
    "to_account": "vendor-mecha-parts",
    "amount_cents": 2500,
}
KEY = "9b6f3c1e-6d3a-4a2f-9f6d-0d0f5f3d2a01"


@pytest.fixture()
def client() -> TestClient:
    app = create_app()
    with TestClient(app) as test_client:
        yield test_client


def test_first_request_executes_and_second_replays(client: TestClient) -> None:
    first = client.post(
        "/api/v1/payments/transfers", json=PAYLOAD, headers={"Idempotency-Key": KEY}
    )
    assert first.status_code == 201
    assert REPLAYED_HEADER not in first.headers

    replay = client.post(
        "/api/v1/payments/transfers", json=PAYLOAD, headers={"Idempotency-Key": KEY}
    )
    assert replay.status_code == 201
    assert replay.headers[REPLAYED_HEADER] == "true"
    assert replay.json() == first.json()
    assert client.app.state.ledger.count() == 1  # exactly-once effect


def test_same_key_different_payload_is_422(client: TestClient) -> None:
    client.post(
        "/api/v1/payments/transfers", json=PAYLOAD, headers={"Idempotency-Key": KEY}
    )
    tampered = dict(PAYLOAD, amount_cents=9999)
    response = client.post(
        "/api/v1/payments/transfers", json=tampered, headers={"Idempotency-Key": KEY}
    )
    assert response.status_code == 422
    assert response.headers["content-type"].startswith("application/problem+json")
    assert client.app.state.ledger.count() == 1


def test_same_key_different_route_is_422(client: TestClient) -> None:
    """The fingerprint covers the route, not just the body."""
    client.post(
        "/api/v1/payments/transfers", json=PAYLOAD, headers={"Idempotency-Key": KEY}
    )
    response = client.post(
        "/api/v1/payments/refunds", json=PAYLOAD, headers={"Idempotency-Key": KEY}
    )
    assert response.status_code == 422


def test_keys_are_scoped_per_account(client: TestClient) -> None:
    """The same key under another account is a different logical request."""
    client.post(
        "/api/v1/payments/transfers",
        json=PAYLOAD,
        headers={"Idempotency-Key": KEY, "X-Account-Id": "acct-alpha"},
    )
    other = client.post(
        "/api/v1/payments/transfers",
        json=PAYLOAD,
        headers={"Idempotency-Key": KEY, "X-Account-Id": "acct-beta"},
    )
    assert other.status_code == 201
    assert REPLAYED_HEADER not in other.headers
    assert client.app.state.ledger.count() == 2


def test_new_key_means_new_operation(client: TestClient) -> None:
    for key in ("key-a", "key-b"):
        response = client.post(
            "/api/v1/payments/transfers", json=PAYLOAD, headers={"Idempotency-Key": key}
        )
        assert response.status_code == 201
    assert client.app.state.ledger.count() == 2


def test_in_flight_key_returns_409(client: TestClient) -> None:
    """Deterministically seed an in-flight row, then collide with it."""
    import json as _json
    import time as _time

    store: IdempotencyStore = client.app.state.store
    body = _json.dumps(PAYLOAD).encode("utf-8")
    fingerprint = fingerprint_request("POST", "/api/v1/payments/transfers", body)
    outcome, _ = store.begin("anonymous", KEY, fingerprint, now=_time.time())
    assert outcome == "started"  # row is now 'in_flight', never completed

    response = client.post(
        "/api/v1/payments/transfers", json=PAYLOAD, headers={"Idempotency-Key": KEY}
    )
    assert response.status_code == 409
    assert response.headers["Retry-After"] == "1"
    assert client.app.state.ledger.count() == 0


def test_expired_key_is_collected_and_reused(client: TestClient) -> None:
    store = IdempotencyStore(ttl_seconds=60)
    app = create_app(store=store)
    with TestClient(app) as expired_client:
        old = expired_client.post(
            "/api/v1/payments/transfers", json=PAYLOAD, headers={"Idempotency-Key": KEY}
        )
        assert old.status_code == 201
        # Simulate time travel: age every row beyond the TTL.
        store._conn.execute("UPDATE idempotency_keys SET created_at = created_at - 120")
        store._conn.commit()
        again = expired_client.post(
            "/api/v1/payments/transfers", json=PAYLOAD, headers={"Idempotency-Key": KEY}
        )
        assert again.status_code == 201
        assert REPLAYED_HEADER not in again.headers  # expired: treated as new
        assert expired_client.app.state.ledger.count() == 2


def test_fifty_concurrent_duplicate_submits_mutate_once(client: TestClient) -> None:
    """50 threads, identical key + payload, released on one barrier.

    Exactly one transfer may exist at the end; every client must observe
    the same receipt body; exactly one response is a non-replayed origin
    response. Duplicates that arrive while the winner is in flight get
    409 and retry until the snapshot is available.
    """
    workers = 50
    barrier = threading.Barrier(workers)
    results: list[tuple[int, bytes, bool]] = []
    results_lock = threading.Lock()

    def submit() -> None:
        barrier.wait()
        attempts = 0
        while True:
            attempts += 1
            assert attempts < 200, "retry storm did not converge"
            response = client.post(
                "/api/v1/payments/transfers",
                json=PAYLOAD,
                headers={"Idempotency-Key": KEY},
            )
            if response.status_code == 409:
                continue  # winner still in flight; poll for the snapshot
            with results_lock:
                results.append(
                    (
                        response.status_code,
                        response.content,
                        response.headers.get(REPLAYED_HEADER) == "true",
                    )
                )
            return

    threads = [threading.Thread(target=submit) for _ in range(workers)]
    for thread in threads:
        thread.start()
    for thread in threads:
        thread.join()

    assert client.app.state.ledger.count() == 1  # one observable mutation
    assert client.app.state.ledger.total_cents() == 2500
    assert len(results) == workers
    statuses = {status for status, _, _ in results}
    bodies = {body for _, body, _ in results}
    assert statuses == {201}
    assert len(bodies) == 1  # every caller received the identical receipt
    non_replayed = sum(1 for _, _, replayed in results if not replayed)
    assert non_replayed == 1  # exactly one origin execution
\end{minted}

\subsection{Listing 10.2 --- Conditional Requests: Strong ETags and the 304 Flow}
\label{sec:ch10-listing2}

A strong validator is a content address: SHA-256 over the exact
representation bytes, rendered from a canonical serializer (sorted keys,
compact separators --- Chapter~4's canonicalization discipline doing double
duty). The demo fetches a part, repeats the fetch conditionally, mutates the
state, and prints the byte ledger: 200 with the full body, 304 with zero
bytes, then a rotated ETag after the write.

\begin{minted}{python}
"""Listing 10.2 -- Conditional requests: strong ETags and the 304 flow.

Northwind Robotics parts-catalog API. The ETag is a strong validator:
SHA-256 over the exact representation bytes, so any byte-level change of
the representation yields a new tag. The demo measures the bandwidth a
304 Not Modified saves on a repeat fetch. Offline, deterministic.
"""

from __future__ import annotations

import hashlib
import json
from dataclasses import dataclass

from fastapi import FastAPI, Request
from fastapi.responses import Response


def strong_etag(representation: bytes) -> str:
    """A strong validator: the quoted SHA-256 of the exact response bytes."""
    return '"' + hashlib.sha256(representation).hexdigest() + '"'


@dataclass
class PartRecord:
    sku: str
    name: str
    unit_price_cents: int
    stock: int

    def render(self) -> bytes:
        """Canonical JSON bytes: one state maps to exactly one byte string."""
        return json.dumps(
            {
                "sku": self.sku,
                "name": self.name,
                "unit_price_cents": self.unit_price_cents,
                "stock": self.stock,
            },
            sort_keys=True,
            separators=(",", ":"),
        ).encode("utf-8")


CATALOG: dict[str, PartRecord] = {
    "NWR-AXLE-42": PartRecord("NWR-AXLE-42", "Heavy-duty drive axle", 189_99, 35),
    "NWR-SERVO-7": PartRecord("NWR-SERVO-7", "Precision servo motor", 74_50, 212),
}


def create_app() -> FastAPI:
    app = FastAPI(title="Northwind Robotics Catalog", version="1.0.0")

    @app.get("/parts/{sku}")
    def get_part(sku: str, request: Request) -> Response:
        part = CATALOG.get(sku)
        if part is None:
            return Response(status_code=404)
        body = part.render()
        etag = strong_etag(body)
        headers = {
            "ETag": etag,
            "Cache-Control": "public, max-age=60",
        }
        if request.headers.get("if-none-match") == etag:
            # Validator matches: the client's copy is current. Send no body.
            return Response(status_code=304, headers=headers)
        return Response(
            content=body,
            media_type="application/json",
            headers=headers,
        )

    @app.post("/parts/{sku}/restock", status_code=200)
    def restock(sku: str, units: int) -> dict[str, object]:
        part = CATALOG[sku]
        part.stock += units  # state change -> new representation -> new ETag
        return {"sku": sku, "stock": part.stock}

    return app


def measure_conditional_savings() -> dict[str, int | float | bool]:
    """Fetch a part twice; the second fetch is a conditional GET.

    Returns byte counts proving the 304 path transfers no representation.
    """
    from fastapi.testclient import TestClient

    snapshot = CATALOG["NWR-AXLE-42"].stock  # keep the module fixture stable
    try:
        with TestClient(create_app()) as client:
            first = client.get("/parts/NWR-AXLE-42")
            assert first.status_code == 200
            etag = first.headers["etag"]
            first_bytes = len(first.content)

            second = client.get(
                "/parts/NWR-AXLE-42", headers={"If-None-Match": etag}
            )
            assert second.status_code == 304
            assert second.headers["etag"] == etag
            second_bytes = len(second.content)

            # The origin state changes: the old validator no longer matches.
            client.post("/parts/NWR-AXLE-42/restock", params={"units": 5})
            third = client.get(
                "/parts/NWR-AXLE-42", headers={"If-None-Match": etag}
            )
            assert third.status_code == 200
            assert third.headers["etag"] != etag

            saved_pct = 100.0 * (1 - second_bytes / first_bytes)
            return {
                "first_status": first.status_code,
                "first_bytes": first_bytes,
                "second_status": second.status_code,
                "second_bytes": second_bytes,
                "saved_pct": round(saved_pct, 1),
                "etag_rotated_after_restock": third.headers["etag"] != etag,
            }
    finally:
        CATALOG["NWR-AXLE-42"].stock = snapshot


if __name__ == "__main__":
    stats = measure_conditional_savings()
    print("Conditional-request bandwidth measurement")
    print("-----------------------------------------")
    print(f"initial GET      : {stats['first_status']} "
          f"({stats['first_bytes']} bytes)")
    print(f"conditional GET  : {stats['second_status']} "
          f"({stats['second_bytes']} bytes)")
    print(f"bandwidth saved  : {stats['saved_pct']}%")
    print(f"ETag rotated on write: {stats['etag_rotated_after_restock']}")
\end{minted}

\begin{minted}{python}
"""Listing 10.2 tests -- ETag stability, 304 flow, rotation on mutation."""

from __future__ import annotations

import pytest
from fastapi.testclient import TestClient

from listing_10_2_conditional_requests import (
    CATALOG,
    create_app,
    measure_conditional_savings,
    strong_etag,
)


@pytest.fixture()
def client() -> TestClient:
    with TestClient(create_app()) as test_client:
        yield test_client
    CATALOG["NWR-AXLE-42"].stock = 35  # restore the seeded fixture


def test_strong_etag_is_content_addressed() -> None:
    assert strong_etag(b"payload") == strong_etag(b"payload")
    assert strong_etag(b"payload") != strong_etag(b"payloae")
    assert strong_etag(b"x").startswith('"') and strong_etag(b"x").endswith('"')


def test_repeat_get_yields_identical_etag(client: TestClient) -> None:
    one = client.get("/parts/NWR-AXLE-42")
    two = client.get("/parts/NWR-AXLE-42")
    assert one.headers["etag"] == two.headers["etag"]
    assert one.headers["cache-control"] == "public, max-age=60"


def test_if_none_match_short_circuits_to_304(client: TestClient) -> None:
    etag = client.get("/parts/NWR-SERVO-7").headers["etag"]
    response = client.get("/parts/NWR-SERVO-7", headers={"If-None-Match": etag})
    assert response.status_code == 304
    assert response.content == b""
    assert response.headers["etag"] == etag  # 304 still carries validators


def test_mutation_rotates_etag_and_invalidates_cache(client: TestClient) -> None:
    before = client.get("/parts/NWR-AXLE-42")
    client.post("/parts/NWR-AXLE-42/restock", params={"units": 10})
    after = client.get(
        "/parts/NWR-AXLE-42", headers={"If-None-Match": before.headers["etag"]}
    )
    assert after.status_code == 200
    assert after.headers["etag"] != before.headers["etag"]
    assert after.json()["stock"] == 45


def test_unknown_part_is_404(client: TestClient) -> None:
    assert client.get("/parts/NOPE-1").status_code == 404


def test_demo_reports_full_savings() -> None:
    stats = measure_conditional_savings()
    assert stats["second_status"] == 304
    assert stats["second_bytes"] == 0
    assert stats["saved_pct"] == 100.0
    assert stats["etag_rotated_after_restock"] is True
\end{minted}

Running the demo prints the measured savings:

\begin{minted}{text}
Conditional-request bandwidth measurement
-----------------------------------------
initial GET      : 200 (88 bytes)
conditional GET  : 304 (0 bytes)
bandwidth saved  : 100.0%
ETag rotated on write: True
\end{minted}

\subsection{Listing 10.3 --- Optimistic Locking with If-Match on Inventory}
\label{sec:ch10-listing3}

The inventory resource stamps every representation with a version-derived
strong ETag (\texttt{"v3"}) and refuses blind writes: no \texttt{If-Match}
is 428, a stale \texttt{If-Match} is 412, and only an exact match reaches
the atomic compare-and-swap. The status codes are not decoration --- each
selects a different client recovery strategy (add the header; re-read and
re-decide; resolve a business conflict).

\begin{minted}{python}
"""Listing 10.3 -- Optimistic locking with If-Match on an inventory API.

Northwind Robotics warehouse inventory. Every representation carries a
version-derived ETag (``"v3"``). Writes must present ``If-Match``:

* header missing           -> 428 Precondition Required,
* header does not match    -> 412 Precondition Failed (lost update refused),
* header matches           -> write applied, version incremented.

The version counter is a strong validator here because the representation
is a pure function of the stored state: equal versions imply byte-identical
bodies. Offline, deterministic.
"""

from __future__ import annotations

import json
import threading
from dataclasses import dataclass

from fastapi import FastAPI, Request
from fastapi.responses import JSONResponse, Response
from pydantic import BaseModel, Field


@dataclass
class InventoryRecord:
    sku: str
    quantity: int
    version: int

    def render(self) -> bytes:
        return json.dumps(
            {"sku": self.sku, "quantity": self.quantity, "version": self.version},
            sort_keys=True,
            separators=(",", ":"),
        ).encode("utf-8")

    @property
    def etag(self) -> str:
        return f'"v{self.version}"'


class InventoryStore:
    """Thread-safe in-memory inventory keyed by SKU."""

    def __init__(self) -> None:
        self._lock = threading.Lock()
        self._records = {
            "NWR-AXLE-42": InventoryRecord("NWR-AXLE-42", quantity=35, version=1),
            "NWR-SERVO-7": InventoryRecord("NWR-SERVO-7", quantity=212, version=1),
        }

    def get(self, sku: str) -> InventoryRecord | None:
        with self._lock:
            record = self._records.get(sku)
            if record is None:
                return None
            return InventoryRecord(record.sku, record.quantity, record.version)

    def compare_and_swap(
        self, sku: str, expected_etag: str, new_quantity: int
    ) -> tuple[str, InventoryRecord | None]:
        """Atomic check-and-set. Returns ('not_found'|'mismatch'|'updated', record)."""
        with self._lock:
            record = self._records.get(sku)
            if record is None:
                return ("not_found", None)
            if record.etag != expected_etag:
                return (
                    "mismatch",
                    InventoryRecord(record.sku, record.quantity, record.version),
                )
            updated = InventoryRecord(sku, new_quantity, record.version + 1)
            self._records[sku] = updated
            return ("updated", updated)


def _problem(status: int, title: str, detail: str, slug: str) -> JSONResponse:
    return JSONResponse(
        status_code=status,
        media_type="application/problem+json",
        content={
            "type": f"https://api.northwind-robotics.example/problems/{slug}",
            "title": title,
            "status": status,
            "detail": detail,
        },
    )


class InventoryUpdate(BaseModel):
    quantity: int = Field(ge=0)


def create_app() -> FastAPI:
    app = FastAPI(title="Northwind Robotics Inventory", version="1.0.0")
    app.state.store = InventoryStore()

    @app.get("/inventory/{sku}")
    def get_inventory(sku: str) -> Response:
        record = app.state.store.get(sku)
        if record is None:
            return Response(status_code=404)
        return Response(
            content=record.render(),
            media_type="application/json",
            headers={"ETag": record.etag},
        )

    @app.put("/inventory/{sku}")
    def put_inventory(
        sku: str, payload: InventoryUpdate, request: Request
    ) -> Response:
        if_match = request.headers.get("if-match")
        if if_match is None:
            return _problem(
                428,
                "Precondition required",
                "Writes to this resource must be conditional: send If-Match"
                " with the ETag of the representation you based your edit on.",
                "precondition-required",
            )
        outcome, record = app.state.store.compare_and_swap(
            sku, if_match, payload.quantity
        )
        if outcome == "not_found":
            return Response(status_code=404)
        if outcome == "mismatch" and record is not None:
            return _problem(
                412,
                "Precondition failed",
                "The resource changed since your copy: current ETag is"
                f" {record.etag}, you sent If-Match: {if_match}. Re-read,"
                " re-decide, and retry.",
                "precondition-failed",
            )
        assert record is not None
        return Response(
            content=record.render(),
            media_type="application/json",
            headers={"ETag": record.etag},
        )

    return app
\end{minted}

\begin{minted}{python}
"""Listing 10.3 tests -- 428/412 selection and lost-update prevention.

The concurrency test drives a deterministic read-modify-write storm:
every writer retries with a fresh read until its compare-and-set wins,
so the final state must equal the sum of all intended increments.
"""

from __future__ import annotations

import threading

import pytest
from fastapi.testclient import TestClient

from listing_10_3_optimistic_locking import create_app

SKU = "NWR-AXLE-42"


@pytest.fixture()
def client() -> TestClient:
    with TestClient(create_app()) as test_client:
        yield test_client


def test_write_without_if_match_is_428(client: TestClient) -> None:
    response = client.put(f"/inventory/{SKU}", json={"quantity": 30})
    assert response.status_code == 428
    assert response.headers["content-type"].startswith("application/problem+json")
    assert client.get(f"/inventory/{SKU}").json()["quantity"] == 35


def test_stale_if_match_is_412_and_state_is_untouched(client: TestClient) -> None:
    etag_v1 = client.get(f"/inventory/{SKU}").headers["etag"]
    # Someone else writes first.
    winner = client.put(
        f"/inventory/{SKU}", json={"quantity": 30}, headers={"If-Match": etag_v1}
    )
    assert winner.status_code == 200
    # Our stale write must be rejected.
    loser = client.put(
        f"/inventory/{SKU}", json={"quantity": 25}, headers={"If-Match": etag_v1}
    )
    assert loser.status_code == 412
    state = client.get(f"/inventory/{SKU}").json()
    assert state["quantity"] == 30  # the winner's value, not the stale one
    assert state["version"] == 2


def test_lost_update_scenario_is_prevented(client: TestClient) -> None:
    """The textbook interleaving: A and B both read v1; A writes; B must lose."""
    read_a = client.get(f"/inventory/{SKU}")
    read_b = client.get(f"/inventory/{SKU}")
    assert read_a.headers["etag"] == read_b.headers["etag"] == '"v1"'

    # A decides from its copy: reserve 5 of 35 -> 30.
    ok = client.put(
        f"/inventory/{SKU}",
        json={"quantity": 35 - 5},
        headers={"If-Match": read_a.headers["etag"]},
    )
    assert ok.status_code == 200
    assert ok.headers["etag"] == '"v2"'

    # B decides from the same stale copy: reserve 8 of 35 -> 27. Refused.
    refused = client.put(
        f"/inventory/{SKU}",
        json={"quantity": 35 - 8},
        headers={"If-Match": read_b.headers["etag"]},
    )
    assert refused.status_code == 412

    # B re-reads, re-decides against current truth, and retries.
    fresh = client.get(f"/inventory/{SKU}")
    retried = client.put(
        f"/inventory/{SKU}",
        json={"quantity": fresh.json()["quantity"] - 8},
        headers={"If-Match": fresh.headers["etag"]},
    )
    assert retried.status_code == 200
    final = client.get(f"/inventory/{SKU}").json()
    assert final["quantity"] == 35 - 5 - 8  # both reservations, none lost
    assert final["version"] == 3


def test_unknown_sku_is_404(client: TestClient) -> None:
    assert client.get("/inventory/NOPE").status_code == 404
    assert (
        client.put(
            "/inventory/NOPE", json={"quantity": 1}, headers={"If-Match": '"v1"'}
        ).status_code
        == 404
    )


def test_negative_quantity_is_rejected(client: TestClient) -> None:
    etag = client.get(f"/inventory/{SKU}").headers["etag"]
    response = client.put(
        f"/inventory/{SKU}", json={"quantity": -1}, headers={"If-Match": etag}
    )
    assert response.status_code == 422


def test_concurrent_read_modify_write_storm_loses_nothing(client: TestClient) -> None:
    """20 threads each increment by one under optimistic retry.

    A writer loops: read -> compute -> conditional write; on 412 it
    re-reads and retries. Loop termination is state-driven (a successful
    compare-and-set), never timing-driven.
    """
    workers = 20
    barrier = threading.Barrier(workers)
    attempts: list[int] = []
    attempts_lock = threading.Lock()

    def increment_once() -> None:
        barrier.wait()
        tries = 0
        while True:
            tries += 1
            assert tries < 500, "optimistic retry loop did not converge"
            current = client.get(f"/inventory/{SKU}")
            etag = current.headers["etag"]
            target = current.json()["quantity"] + 1
            written = client.put(
                f"/inventory/{SKU}",
                json={"quantity": target},
                headers={"If-Match": etag},
            )
            if written.status_code == 200:
                with attempts_lock:
                    attempts.append(tries)
                return
            assert written.status_code == 412  # only contention may fail us

    threads = [threading.Thread(target=increment_once) for _ in range(workers)]
    for thread in threads:
        thread.start()
    for thread in threads:
        thread.join()

    final = client.get(f"/inventory/{SKU}").json()
    assert final["quantity"] == 35 + workers  # every increment landed
    assert final["version"] == 1 + workers    # one version per accepted write
    assert len(attempts) == workers
\end{minted}

\subsection{Listing 10.4 --- Double-Submit Storm Harness}
\label{sec:ch10-listing4}

This harness reproduces the incident shape --- fifty identical
\texttt{POST /transfers} requests fired from one barrier --- twice: once as
the unsafe legacy client (no key, every duplicate executes) and once with
the \texttt{Idempotency-Key}. The output is the chapter's thesis in two
lines of measured output.

\begin{minted}{python}
"""Listing 10.4 -- Double-submit storm harness.

Reproduces the classic checkout bug -- a user double-clicks "Pay", the
load balancer retries a timed-out request, or a mobile client re-sends on
flaky Wi-Fi -- and proves that an Idempotency-Key turns a storm of N
duplicate requests into exactly one observable mutation.

The harness is deterministic: workers are released on one
threading.Barrier, retries on 409 are state-driven (poll until the winner
publishes its snapshot), and every assertion is on final state, never on
timing. Run directly for a printed before/after comparison.
"""

from __future__ import annotations

import threading
from dataclasses import dataclass, field

from fastapi.testclient import TestClient

from listing_10_1_idempotency_middleware import REPLAYED_HEADER, create_app

TRANSFERS_URL = "/api/v1/payments/transfers"
PAYLOAD = {
    "from_account": "ops-checking",
    "to_account": "vendor-mecha-parts",
    "amount_cents": 2500,
}


@dataclass(frozen=True)
class StormResult:
    """What a storm of duplicate submits produced."""

    workers: int
    status_codes: list[int] = field(default_factory=list)
    bodies: set[bytes] = field(default_factory=set)
    replayed_count: int = 0
    ledger_entries: int = 0
    ledger_total_cents: int = 0


def run_double_submit_storm(
    workers: int = 50, use_idempotency_key: bool = True
) -> StormResult:
    """Fire ``workers`` identical POSTs at the API from a single barrier.

    With ``use_idempotency_key=False`` this simulates the unsafe legacy
    client: every request executes. With the key, exactly one executes.
    """
    app = create_app()
    key = "storm-" + ("keyed" if use_idempotency_key else "naive")

    with TestClient(app) as client:
        barrier = threading.Barrier(workers)
        status_codes: list[int] = []
        bodies: set[bytes] = set()
        replayed = 0
        lock = threading.Lock()

        def submit() -> None:
            nonlocal replayed
            headers = (
                {"Idempotency-Key": key} if use_idempotency_key else {}
            )
            barrier.wait()
            attempts = 0
            while True:
                attempts += 1
                assert attempts < 200, "storm did not converge"
                response = client.post(TRANSFERS_URL, json=PAYLOAD, headers=headers)
                if response.status_code == 409 and use_idempotency_key:
                    continue  # winner in flight; poll for its snapshot
                with lock:
                    status_codes.append(response.status_code)
                    bodies.add(response.content)
                    if response.headers.get(REPLAYED_HEADER) == "true":
                        replayed += 1
                return

        threads = [threading.Thread(target=submit) for _ in range(workers)]
        for thread in threads:
            thread.start()
        for thread in threads:
            thread.join()

        return StormResult(
            workers=workers,
            status_codes=status_codes,
            bodies=bodies,
            replayed_count=replayed,
            ledger_entries=app.state.ledger.count(),
            ledger_total_cents=app.state.ledger.total_cents(),
        )


def main() -> None:
    print("Double-submit storm: 50 identical POST /transfers")
    print("=" * 52)
    naive = run_double_submit_storm(workers=50, use_idempotency_key=False)
    print(f"without Idempotency-Key : {naive.ledger_entries} ledger entries,"
          f" {naive.ledger_total_cents} cents moved  <-- the incident")
    keyed = run_double_submit_storm(workers=50, use_idempotency_key=True)
    print(f"with    Idempotency-Key : {keyed.ledger_entries} ledger entry,"
          f" {keyed.ledger_total_cents} cents moved, "
          f"{keyed.replayed_count} replayed responses")


if __name__ == "__main__":
    main()
\end{minted}

\begin{minted}{python}
"""Listing 10.4 tests -- the storm harness proves effectively-once."""

from __future__ import annotations

from listing_10_4_double_submit_storm import run_double_submit_storm


def test_naive_client_executes_every_duplicate() -> None:
    """The unsafe baseline: no key, so all 50 duplicates charge through."""
    result = run_double_submit_storm(workers=50, use_idempotency_key=False)
    assert result.ledger_entries == 50
    assert result.ledger_total_cents == 50 * 2500
    assert set(result.status_codes) == {201}
    assert result.replayed_count == 0


def test_keyed_storm_mutates_exactly_once() -> None:
    result = run_double_submit_storm(workers=50, use_idempotency_key=True)
    assert result.ledger_entries == 1
    assert result.ledger_total_cents == 2500
    assert set(result.status_codes) == {201}
    assert len(result.bodies) == 1  # identical receipt for every caller
    # Exactly one origin execution; the other 49 observed the replay.
    assert result.replayed_count == 49


def test_small_storms_also_collapse_to_one_mutation() -> None:
    for workers in (2, 3, 7):
        result = run_double_submit_storm(workers=workers, use_idempotency_key=True)
        assert result.ledger_entries == 1
        assert result.replayed_count == workers - 1
\end{minted}

\begin{minted}{text}
Double-submit storm: 50 identical POST /transfers
====================================================
without Idempotency-Key : 50 ledger entries, 125000 cents moved  <-- the incident
with    Idempotency-Key : 1 ledger entry, 2500 cents moved, 49 replayed responses
\end{minted}

\subsection{Listing 10.5 --- Cache-Policy Simulator}
\label{sec:ch10-listing5}

The simulator makes RFC~9111 executable: parse a \texttt{Cache-Control}
policy, script an origin's behavior over a logical timeline (representations
change at $t=70$, the origin falls over at $t=130$), and compute the exact
outcome of every request --- HIT, MISS, revalidation hit or miss,
stale-while-revalidate, stale-if-error, bypass. Because time is a logical
integer and the origin is a sorted script, every run is bit-for-bit
reproducible, which makes the simulator a teaching instrument and a policy
test harness at once.

\begin{minted}{python}
"""Listing 10.5 -- Cache-policy simulator (pure stdlib).

Given a response's cache policy (Cache-Control directives + validators)
and a deterministic request timeline against a scripted origin, compute
the outcome of every request: HIT, MISS, revalidation, stale serve, or
bypass. Teaches RFC 9111 freshness semantics without a network.

Time is a logical integer (seconds); the origin is a script of
(t, state) changes, so every run is bit-for-bit reproducible.
"""

from __future__ import annotations

from collections import Counter
from dataclasses import dataclass
from typing import Optional

HEURISTIC_FRACTION = 0.1  # RFC 9111 heuristic: 10% of (Date - Last-Modified)


# --------------------------------------------------------------------------
# Policy parsing
# --------------------------------------------------------------------------
@dataclass(frozen=True)
class CachePolicy:
    """The cache-relevant slice of one origin response."""

    no_store: bool = False
    no_cache: bool = False
    must_revalidate: bool = False
    public: bool = False
    private: bool = False
    max_age: Optional[float] = None
    stale_while_revalidate: Optional[float] = None
    stale_if_error: Optional[float] = None
    etag: Optional[str] = None
    last_modified: Optional[float] = None
    date: float = 0.0

    def freshness_lifetime(self) -> float:
        """Explicit max-age wins; otherwise the RFC 9111 heuristic."""
        if self.max_age is not None:
            return self.max_age
        if self.last_modified is not None:
            return HEURISTIC_FRACTION * max(0.0, self.date - self.last_modified)
        return 0.0


def parse_cache_control(value: str) -> dict[str, Optional[str]]:
    """'public, max-age=60' -> {'public': None, 'max-age': '60'}."""
    directives: dict[str, Optional[str]] = {}
    for token in value.split(","):
        name, _, arg = token.strip().partition("=")
        if name:
            directives[name.lower()] = arg.strip('"') if arg else None
    return directives


def policy_from_headers(
    cache_control: str,
    *,
    etag: Optional[str] = None,
    last_modified: Optional[float] = None,
    date: float = 0.0,
) -> CachePolicy:
    d = parse_cache_control(cache_control)
    return CachePolicy(
        no_store="no-store" in d,
        no_cache="no-cache" in d,
        must_revalidate="must-revalidate" in d,
        public="public" in d,
        private="private" in d,
        max_age=float(d["max-age"]) if d.get("max-age") is not None else None,
        stale_while_revalidate=(
            float(d["stale-while-revalidate"])
            if d.get("stale-while-revalidate") is not None
            else None
        ),
        stale_if_error=(
            float(d["stale-if-error"])
            if d.get("stale-if-error") is not None
            else None
        ),
        etag=etag,
        last_modified=last_modified,
        date=date,
    )


# --------------------------------------------------------------------------
# Scripted origin
# --------------------------------------------------------------------------
@dataclass(frozen=True)
class OriginState:
    status: int  # 200 or 500
    etag: str
    body_version: int


class ScriptedOrigin:
    """A deterministic origin: a sorted script of (t, OriginState) changes."""

    def __init__(self, script: list[tuple[float, OriginState]]) -> None:
        if not script:
            raise ValueError("origin script needs at least one state")
        self._script = sorted(script, key=lambda item: item[0])
        self.fetch_count = 0

    def state_at(self, t: float) -> OriginState:
        state = self._script[0][1]
        for change_at, candidate in self._script:
            if change_at <= t:
                state = candidate
            else:
                break
        return state

    def fetch(self, t: float) -> OriginState:
        self.fetch_count += 1
        return self.state_at(t)


# --------------------------------------------------------------------------
# The cache
# --------------------------------------------------------------------------
@dataclass
class StoredEntry:
    policy: CachePolicy
    etag: str
    body_version: int
    stored_at: float


HIT = "HIT"
MISS = "MISS"
REVALIDATED_HIT = "REVALIDATED_HIT"    # conditional GET -> 304
REVALIDATED_MISS = "REVALIDATED_MISS"  # conditional GET -> 200 (new ETag)
STALE_WHILE_REVALIDATE = "STALE_WHILE_REVALIDATE"
STALE_IF_ERROR = "STALE_IF_ERROR"
ERROR = "ERROR"     # origin down, nothing servable
BYPASS = "BYPASS"   # no-store: never reused


class SimulatedCache:
    """Single-entry shared cache executing RFC 9111 freshness decisions."""

    def __init__(self, origin: ScriptedOrigin) -> None:
        self.origin = origin
        self.stored: Optional[StoredEntry] = None
        self.stats: Counter[str] = Counter()

    def request(self, t: float) -> str:
        outcome = self._serve(t)
        self.stats[outcome] += 1
        return outcome

    def _serve(self, t: float) -> str:
        entry = self.stored
        if entry is None or entry.policy.no_store:
            return self._fetch_and_store(t, miss_label=MISS if entry is None else BYPASS)

        policy = entry.policy
        age = t - entry.stored_at
        lifetime = policy.freshness_lifetime()
        fresh = age <= lifetime and not policy.no_cache

        if fresh:
            return HIT

        # Stale-while-revalidate: serve stale now, refresh in the background.
        if (
            policy.stale_while_revalidate is not None
            and not policy.no_cache
            and not policy.must_revalidate
            and age <= lifetime + policy.stale_while_revalidate
        ):
            self._background_revalidate(t, entry)
            return STALE_WHILE_REVALIDATE

        return self._revalidate(t, entry, age, lifetime)

    def _revalidate(
        self, t: float, entry: StoredEntry, age: float, lifetime: float
    ) -> str:
        policy = entry.policy
        origin_state = self.origin.fetch(t)
        if origin_state.status >= 500:
            if (
                policy.stale_if_error is not None
                and not policy.must_revalidate
                and age <= lifetime + policy.stale_if_error
            ):
                return STALE_IF_ERROR
            return ERROR
        if origin_state.etag == entry.etag:
            # 304: freshen the stored entry in place.
            entry.stored_at = t
            return REVALIDATED_HIT
        # 200: a new representation replaces the entry.
        self.stored = StoredEntry(
            policy=entry.policy,
            etag=origin_state.etag,
            body_version=origin_state.body_version,
            stored_at=t,
        )
        return REVALIDATED_MISS

    def _background_revalidate(self, t: float, entry: StoredEntry) -> None:
        """Model the async refresh synchronously for determinism."""
        origin_state = self.origin.fetch(t)
        if origin_state.status < 500:
            self.stored = StoredEntry(
                policy=entry.policy,
                etag=origin_state.etag,
                body_version=origin_state.body_version,
                stored_at=t,
            )

    def _fetch_and_store(self, t: float, miss_label: str) -> str:
        origin_state = self.origin.fetch(t)
        if origin_state.status >= 500:
            return ERROR
        policy = CachePolicy(
            max_age=60, etag=origin_state.etag, date=t
        )  # demo default; scenarios override via inject()
        if self.stored is not None and self.stored.policy.no_store:
            return miss_label  # BYPASS: policy forbids storage entirely
        self.stored = StoredEntry(
            policy=policy,
            etag=origin_state.etag,
            body_version=origin_state.body_version,
            stored_at=t,
        )
        return miss_label

    def inject(self, policy: CachePolicy, t: float) -> None:
        """Prime the cache with an entry under an explicit policy."""
        state = self.origin.state_at(t)
        self.stored = StoredEntry(
            policy=policy,
            etag=state.etag,
            body_version=state.body_version,
            stored_at=t,
        )


def run_timeline(cache: SimulatedCache, times: list[float]) -> list[str]:
    return [cache.request(t) for t in times]


if __name__ == "__main__":
    origin = ScriptedOrigin(
        [
            (0, OriginState(200, '"v1"', 1)),
            (70, OriginState(200, '"v2"', 2)),
            (130, OriginState(500, '"v2"', 2)),
        ]
    )
    cache = SimulatedCache(origin)
    cache.request(0)  # MISS, stores v1 with max-age=60
    print("t=10  :", cache.request(10))    # HIT
    print("t=80  :", cache.request(80))    # revalidate -> new etag v2
    print("t=150 :", cache.request(150))   # stale, origin down
    print("stats :", dict(cache.stats))
\end{minted}

\begin{minted}{python}
"""Listing 10.5 tests -- directive semantics, proven on scripted timelines."""

from __future__ import annotations

from listing_10_5_cache_simulator import (
    BYPASS,
    ERROR,
    HIT,
    MISS,
    OriginState,
    REVALIDATED_HIT,
    REVALIDATED_MISS,
    STALE_IF_ERROR,
    STALE_WHILE_REVALIDATE,
    SimulatedCache,
    ScriptedOrigin,
    parse_cache_control,
    policy_from_headers,
    run_timeline,
)


def make_cache(
    script: list[tuple[float, OriginState]] | None = None,
) -> SimulatedCache:
    if script is None:
        script = [(0, OriginState(200, '"v1"', 1))]
    return SimulatedCache(ScriptedOrigin(script))


def test_parse_cache_control() -> None:
    d = parse_cache_control('public, max-age=60, stale-while-revalidate=30')
    assert d == {"public": None, "max-age": "60", "stale-while-revalidate": "30"}


def test_max_age_governs_freshness() -> None:
    cache = make_cache()
    cache.inject(policy_from_headers("public, max-age=60", etag='"v1"'), t=0)
    assert run_timeline(cache, [10, 60]) == [HIT, HIT]
    assert cache.request(61) == REVALIDATED_HIT  # stale -> conditional GET -> 304
    assert cache.request(70) == HIT  # freshened by the 304


def test_no_store_never_caches() -> None:
    cache = make_cache()
    cache.inject(policy_from_headers("no-store", etag='"v1"'), t=0)
    assert run_timeline(cache, [1, 2, 3]) == [BYPASS, BYPASS, BYPASS]
    assert cache.origin.fetch_count == 3  # every request hits the origin


def test_no_cache_revalidates_every_time() -> None:
    cache = make_cache()
    cache.inject(
        policy_from_headers("no-cache", etag='"v1"', date=0), t=0
    )
    # Validators make each revalidation cheap (304), but nothing is served
    # without contacting the origin first.
    assert run_timeline(cache, [1, 2]) == [REVALIDATED_HIT, REVALIDATED_HIT]
    assert cache.origin.fetch_count == 2


def test_heuristic_freshness_without_max_age() -> None:
    cache = make_cache()
    # Last-Modified 1000s before Date -> heuristic freshness = 100s.
    policy = policy_from_headers(
        "public", etag='"v1"', last_modified=0.0, date=1000.0
    )
    cache.inject(policy, t=1000.0)
    assert cache.request(1050.0) == HIT
    assert cache.request(1200.0) == REVALIDATED_HIT


def test_changed_etag_yields_revalidated_miss() -> None:
    cache = make_cache(
        [
            (0, OriginState(200, '"v1"', 1)),
            (50, OriginState(200, '"v2"', 2)),
        ]
    )
    cache.inject(policy_from_headers("public, max-age=10", etag='"v1"'), t=0)
    assert cache.request(60) == REVALIDATED_MISS  # 200 with new body
    assert cache.stored is not None and cache.stored.etag == '"v2"'
    assert cache.request(65) == HIT  # the fresh replacement now serves


def test_stale_while_revalidate_serves_stale_then_refreshes() -> None:
    cache = make_cache(
        [
            (0, OriginState(200, '"v1"', 1)),
            (40, OriginState(200, '"v2"', 2)),
        ]
    )
    cache.inject(
        policy_from_headers(
            "public, max-age=30, stale-while-revalidate=60", etag='"v1"'
        ),
        t=0,
    )
    assert cache.request(50) == STALE_WHILE_REVALIDATE  # stale served...
    assert cache.stored is not None and cache.stored.etag == '"v2"'  # ...refreshed
    assert cache.request(55) == HIT


def test_stale_if_error_rescues_an_origin_outage() -> None:
    cache = make_cache(
        [
            (0, OriginState(200, '"v1"', 1)),
            (100, OriginState(500, '"v1"', 1)),
        ]
    )
    cache.inject(
        policy_from_headers("public, max-age=30, stale-if-error=300", etag='"v1"'),
        t=0,
    )
    assert cache.request(120) == STALE_IF_ERROR  # origin down, stale saved us


def test_must_revalidate_forbids_stale_if_error() -> None:
    cache = make_cache(
        [
            (0, OriginState(200, '"v1"', 1)),
            (100, OriginState(500, '"v1"', 1)),
        ]
    )
    cache.inject(
        policy_from_headers(
            "public, max-age=30, must-revalidate, stale-if-error=300", etag='"v1"'
        ),
        t=0,
    )
    assert cache.request(120) == ERROR  # must-revalidate wins: no stale serve


def test_full_demo_timeline_is_deterministic() -> None:
    def run() -> tuple[list[str], dict[str, int]]:
        origin = ScriptedOrigin(
            [
                (0, OriginState(200, '"v1"', 1)),
                (70, OriginState(200, '"v2"', 2)),
                (130, OriginState(500, '"v2"', 2)),
            ]
        )
        cache = SimulatedCache(origin)
        outcomes = run_timeline(cache, [0, 10, 80, 100, 150])
        return outcomes, dict(cache.stats)

    first_outcomes, first_stats = run()
    second_outcomes, second_stats = run()
    assert first_outcomes == second_outcomes  # bit-for-bit reproducible
    assert first_stats == second_stats
    assert first_outcomes == [MISS, HIT, REVALIDATED_MISS, HIT, ERROR]
\end{minted}

\begin{keyconceptbox}
All five listings share one architectural motif: \emph{correctness lives in
an atomic compare-and-set at a storage boundary}. The idempotency store
claims keys with a primary-key \texttt{INSERT}; the inventory store swaps
state under one lock with the expected version in hand; the cache revalidates
against a validator rather than trusting a clock. Distribution changes the
storage, not the shape of the argument.
\end{keyconceptbox}

%% ----------------------------------------------------------------------------
\section{Common Pitfalls \& Anti-Patterns}
\label{sec:ch10-pitfalls}

\subsection{Storing the Response Before Committing the Mutation}
Writing the idempotency snapshot before the business transaction commits
inverts the failure window: a crash now leaves a stored success response for
a mutation that rolled back. The retry then replays a receipt for a transfer
that never happened. Rule: claim the key early (\texttt{in\_flight}), record
the response only with (or after) the durable commit, and abandon the claim
on failure so the client can retry cleanly.

\subsection{Reusing Idempotency Keys Across Endpoints}
A key must identify one logical operation on one endpoint. Reusing a key on
a second endpoint either replays the wrong response (if fingerprints are not
checked) or 422s (if they are). Clients should mint keys per
\emph{user intent} --- one per checkout form instance --- and servers must
enforce the fingerprint so sloppy reuse is loud, not silent.

\subsection{Hashing Only the Body, Not the Route}
A fingerprint computed from the request body alone lets the same key +
identical JSON body collide across \texttt{/transfers} and
\texttt{/refunds}. Listing~10.1's \texttt{fingerprint\_request} hashes
method and path alongside the canonical body for exactly this reason; the
test \texttt{test\_same\_key\_different\_route\_is\_422} is the regression
gate.

\subsection{Retrying POST Without Keys}
An HTTP client library that transparently retries POST on timeout is an
undocumented double-charge generator. Audit every layer --- your SDK, your
proxy, your load balancer's retry policy --- and either mint an idempotency
key per logical operation or disable retries on unsafe methods. Retry
budgets (Chapter~2) limit \emph{how often}; keys limit \emph{what repeats
mean}.

\subsection{Weak ETags on Byte-Sensitive Resources}
A weak validator (\texttt{W/"abc"}) asserts semantic equivalence, not byte
equality. Using one on a resource that feeds range requests, delta
encodings, or cryptographic verification lets a cache combine stale bytes
with fresh offsets. Rule: content-addressed (hash) or version-addressed
ETags for anything byte-sensitive; weak ETags only for representations that
legitimately vary in bytes while staying semantically equal.

\subsection{The no-store / no-cache Confusion}
\texttt{no-cache} \emph{stores} the response and revalidates on every use;
\texttt{no-store} never stores. Teams that write \texttt{no-cache} on
responses containing personal data have built a compliant-looking privacy
leak: the bytes sit in every intermediary on the path. When the data is
sensitive, the directive is \texttt{no-store}, full stop; add
\texttt{private} when per-user browser caching is acceptable.

\subsection{Client-Generated ETags}
Validators are issued by the authority that owns the state --- the server.
A client that invents an ETag (hashing its own serialization, guessing a
version) couples the cache protocol to a serialization the server never
committed to, and optimistic locking degenerates into a coin flip. Clients
treat ETags as opaque; servers mint them; nobody else touches them.

\subsection{Distributed Locks Without Fencing Tokens}
A lease-based lock makes a paused holder and a crashed holder
indistinguishable; when the lease lapses, two clients believe they hold the
lock. Without fencing tokens checked at the resource, the second acquisition
does not restore mutual exclusion --- it merely narrows the race window.
If the resource cannot check tokens, the lock is advisory and must be
treated as a throughput optimization, never as a correctness
mechanism~\cite{kleppmann2017ddia}.

%% ----------------------------------------------------------------------------
\section{Hands-On Production Lab \& Real-World Case Study}
\label{sec:ch10-lab}

\subsection{Lab: An Idempotent Payments API That Survives Retry Storms}

\textbf{Scenario.} Northwind Robotics exposes
\texttt{POST /api/v1/payments/transfers} to its dealer network. Dealers run
flaky warehouse Wi-Fi; their integration library retries aggressively. Your
job: ship an endpoint on which \emph{any} number of duplicate submissions
produces exactly one transfer.

\textbf{Build order.}
\begin{enumerate}
  \item Start from Listing~10.1. Run
        \texttt{pytest test\_listing\_10\_1.py -q} and confirm all eight
        tests pass.
  \item Extend \texttt{TransferRequest} with a \texttt{reference} field
        (dealer-supplied order number) and prove by test that canonical JSON
        fingerprinting treats reordered keys as the same payload.
  \item Lower the TTL to 5 minutes and write a test (via the
        \texttt{now} parameter, not \texttt{sleep}) showing an aged-out key
        is collected and the operation executes again.
  \item Add metrics: count executions, replays, mismatches, and in-flight
        rejections; expose them at \texttt{GET /api/v1/\_meta/idempotency}.
  \item Run Listing~10.4's storm at 200 workers against your extended app.
        Acceptance: one ledger entry, 200 identical receipts.
\end{enumerate}

\textbf{Deliverable.} The extended middleware, the metrics endpoint, and the
passing storm output pasted into your lab notes.

\subsection{War Story: The Double-Charged Customers}
\label{sec:ch10-warstory}

A payments team (the details are composited from several public incident
reviews) shipped a new load balancer configured to retry requests on
connection reset --- a well-intentioned resilience default. Retries were
enabled for \emph{all} methods, including POST, on the theory that ``the
API is mostly idempotent anyway.'' The charges endpoint was not. For three
days, a small fraction of checkouts --- those whose responses coincided
with a connection reset --- executed twice. Because both executions
returned 201 with distinct charge IDs, no client-side error ever fired; the
first symptom was customers reporting duplicate card charges, hours later,
through the support queue. Forensics was hampered because the two charges
had different IDs and timestamps seconds apart: they looked like two
deliberate purchases. Remediation cost: a weekend of ledger forensics,
thousands of manual refunds, and a lasting support burden.

Three structural fixes followed. (1)~Retries at every infrastructure layer
were restricted to idempotent methods unless the request carries an
idempotency key. (2)~The charges endpoint adopted idempotency keys with a
24-hour TTL --- the client SDK mints one per checkout session. (3)~Replay
telemetry (\texttt{Idempotency-Replayed} counts) became a dashboard: a
replay spike is now an early-warning signal for exactly the network
conditions that used to produce double charges. The lesson generalizes:
\emph{infrastructure retries are invisible application behavior}, and any
retry you did not design is a retry you have not made safe.

%% ----------------------------------------------------------------------------
\section{Self-Assessment Exercises \& Deep-Dive Questions}
\label{sec:ch10-exercises}

\begin{enumerate}
  \item \textbf{[junior]} Classify each as safe, idempotent, both, or
        neither: GET; PUT; DELETE; POST; \texttt{PATCH} with body
        \texttt{\{"op": "set", "field": "qty", "value": 5\}};
        \texttt{PATCH} with body \texttt{\{"op": "inc", "field": "qty"\}}.
  \item \textbf{[junior]} Your endpoint returns
        \texttt{Cache-Control: no-cache}. A colleague claims the response
        is never stored anywhere. Correct them precisely.
  \item \textbf{[junior]} Given \texttt{max-age=60} and a response stored
        at $t=0$, which of the requests at $t \in \{30, 60, 61\}$ hit the
        cache? Verify with Listing~10.5.
  \item \textbf{[mid]} A client reuses one idempotency key for its whole
        batch job, varying only the payload. What happens, request by
        request, against Listing~10.1? Which status codes appear and why?
  \item \textbf{[mid]} Listing~10.1 abandons the key row on downstream
        5xx. Argue both sides: should a deterministic 422 (validation
        failure) be stored and replayed, or abandoned? What does Stripe's
        design choose, and why?
  \item \textbf{[mid]} Extend Listing~10.5 with a \texttt{Vary} model: the
        cache key becomes (URI, Accept-Language). Show a timeline in which
        omitting \texttt{Vary} serves a French representation to an English
        requester.
  \item \textbf{[senior]} Design the database schema for idempotency keys
        when the mutation and the snapshot must commit atomically in one
        relational transaction. Where does the \texttt{in\_flight} row
        live, and what does a crashed transaction leave behind?
  \item \textbf{[senior]} A hot inventory SKU suffers livelock under
        optimistic retries: 95\% of conditional writes return 412. Compare
        three escapes --- bounded retry with jitter, pessimistic row locks,
        and per-SKU command queues --- and pick one with numbers.
  \item \textbf{[staff]} Your CDN supports surrogate keys. Design the
        tagging and purge scheme for the Northwind catalog so that a price
        change on one part invalidates exactly the affected URLs across
        product, search, and recommendation surfaces.
  \item \textbf{[staff]} Prove or refute: ``If every mutation endpoint
        requires an idempotency key, we no longer need unique constraints
        in the database.'' Treat the key TTL as the crux of your argument.
\end{enumerate}

%% ----------------------------------------------------------------------------
\section{Interview Q\&A Vault}
\label{sec:ch10-interview}

\subsection*{Junior (Q1--Q10)}
\begin{enumerate}[label=\textbf{Q\arabic*.}, leftmargin=1.6cm]
  \item \textbf{Define idempotency in one sentence.}
        Repeating an operation any number of times leaves the observable
        state exactly where one application left it: $f(f(x)) = f(x)$.
  \item \textbf{Which HTTP methods are idempotent by specification?}
        GET, HEAD, OPTIONS, TRACE (also safe), plus PUT and DELETE. POST is
        not; PATCH is not in general~\cite{rfc9110}.
  \item \textbf{Prove POST is not idempotent by spec.}
        Exhibit a counterexample permitted by the spec:
        \texttt{POST /orders} creates a new order per request. Two
        identical requests produce two orders with distinct IDs ---
        $f(f(x)) \neq f(x)$ on the observable state (order count). Since
        the specification assigns POST no fixed-point guarantee, no
        conforming server is required to deduplicate, QED.
  \item \textbf{Safe versus idempotent --- what is the difference?}
        Safe methods change nothing (zero applications $\equiv$ one).
        Idempotent methods may change state once but not additionally on
        repetition. GET is both; PUT is idempotent but not safe.
  \item \textbf{Is a second DELETE returning 404 still idempotent?}
        Yes. Idempotency constrains the \emph{state}, not the status code:
        after the first DELETE the resource is absent, and it stays absent
        no matter how often the request repeats.
  \item \textbf{What is an idempotency key and who generates it?}
        A unique token (typically UUIDv4) generated by the \emph{client},
        once per logical operation, sent on every retry of that operation
        so the server can recognize and collapse duplicates.
  \item \textbf{What does \texttt{Cache-Control: no-store} mean?}
        Caches must not store any part of the response; every request must
        reach the origin. It is the correct directive for credentials and
        personal data.
  \item \textbf{What does a 304 Not Modified response contain?}
        No body --- only headers, including refreshed validators and cache
        directives. It tells the cache its stored representation is still
        current.
  \item \textbf{What is an ETag?}
        An opaque, server-issued validator identifying one version of a
        representation. A strong ETag changes when a single byte changes;
        a weak one (\texttt{W/} prefix) asserts only semantic equivalence.
  \item \textbf{What status code answers a failed \texttt{If-Match}?}
        412 Precondition Failed: the resource changed since the client's
        copy; the client should re-read, re-decide, and retry.
\end{enumerate}

\subsection*{Mid (Q11--Q22)}
\begin{enumerate}[label=\textbf{Q\arabic*.}, leftmargin=1.6cm, start=11]
  \item \textbf{\texttt{no-cache} versus \texttt{no-store}?}
        \texttt{no-cache} permits storage but requires revalidation before
        every reuse; \texttt{no-store} forbids storage entirely. The names
        are famously misleading --- memorize the semantics, not the words.
  \item \textbf{Design idempotency for a payment endpoint.}
        Client mints a UUID per checkout intent and sends it as
        \texttt{Idempotency-Key}. Server: atomic check-and-insert of
        (account, key) with a fingerprint of method+route+canonical body;
        execute once; snapshot status+body in the same transaction; replay
        on duplicate; 422 on fingerprint mismatch; 409 while the twin is in
        flight; TTL of 24h matched to the client's retry horizon.
  \item \textbf{412 versus 409 versus 428?}
        428: the server requires conditional requests and none was sent.
        412: a precondition (\texttt{If-Match}) was sent and evaluated
        false. 409: the request conflicts with server \emph{state} ---
        business-rule conflicts and in-flight idempotency collisions.
        Precondition problems are 412/428; state conflicts are 409.
  \item \textbf{Why is exactly-once delivery impossible?}
        The Two Generals' Problem: over an unreliable channel, no finite
        message exchange lets both parties be certain the other will act.
        After a timeout the sender cannot distinguish ``request lost'' from
        ``response lost,'' so no protocol guarantees exactly-once. Systems
        ship at-least-once delivery plus idempotent processing:
        effectively-once~\cite{kleppmann2017ddia}.
  \item \textbf{What are \texttt{stale-while-revalidate} semantics?}
        Within the directive's window past expiry, the cache serves the
        stale representation \emph{immediately} and refreshes it from the
        origin in the background. Latency stays at cache speed; staleness
        is bounded by the window; the next request gets the fresh copy.
  \item \textbf{What is \texttt{stale-if-error} for?}
        Availability during origin outages: if revalidation fails with an
        error, the cache may serve the stale entry within the window.
        \texttt{must-revalidate} disables this escape hatch --- choose it
        when a stale answer is worse than an error (e.g.\ account
        balances).
  \item \textbf{Strong versus weak ETags --- when does the difference
        matter?}
        Range requests and byte-level combining: a weak validator cannot
        prove two byte strings identical, so caches must not use it to
        satisfy partial-content requests. Content-addressed (hash) ETags
        are strong by construction.
  \item \textbf{How would you fingerprint an idempotent request?}
        SHA-256 over method, route template, and the canonicalized body
        (JSON with sorted keys and fixed separators). Hashing the body
        alone misses cross-endpoint collisions; hashing raw bytes misses
        semantically identical JSON with reordered keys.
  \item \textbf{What is the lost-update problem?}
        Two transactions read the same state, compute new states privately,
        and write; the second write overwrites the first, silently
        discarding one update. Prevented by conditional writes
        (compare-and-set on a version) or by serializing access with locks.
  \item \textbf{What is a fencing token?}
        A monotonic token issued at lock acquisition and checked by the
        resource on every write; writes bearing an older token are
        rejected. It closes the paused-lock-holder hole that lease expiry
        alone cannot~\cite{kleppmann2017ddia}.
  \item \textbf{Your client retries timed-out POSTs. What can go wrong and
        how do you fix it?}
        A timeout means the outcome is unknown; a blind retry can apply the
        operation twice (double charge). Fix: mint an idempotency key per
        logical operation so the server collapses the retry into a replay,
        or move the operation behind a naturally idempotent PUT.
  \item \textbf{What is cache poisoning via unkeyed input?}
        When a response depends on data not present in the cache key or the
        \texttt{Vary} set (a reflected header, an ignored query
        parameter), an attacker can store a malicious representation under
        an innocent key and have the cache serve it to other users.
        Defense: every response byte is a function of key plus
        \texttt{Vary}; never reflect unkeyed input.
\end{enumerate}

\subsection*{Senior (Q23--Q28)}
\begin{enumerate}[label=\textbf{Q\arabic*.}, leftmargin=1.6cm, start=23]
  \item \textbf{Why must the idempotency snapshot and the business mutation
        commit in one transaction?}
        Otherwise a crash between the two leaves an unrecoverable skew: a
        replayable receipt for a rolled-back mutation, or a mutation with
        no receipt whose \texttt{in\_flight} key blocks retries forever.
        Atomicity converts the pair into one fact~\cite{kleppmann2017ddia}.
  \item \textbf{Where should the dedup check live: middleware, handler, or
        database?}
        The database owns correctness (the unique constraint is the atomic
        compare-and-set); middleware owns the protocol mapping (422/409/
        replay); the handler should stay ignorant. Pushing the check into
        application code above the store reintroduces TOCTOU races.
  \item \textbf{How do you size the idempotency-key TTL?}
        From the client's worst-case retry horizon: maximum backoff
        schedule times maximum attempts, plus clock skew, plus a margin
        for human-in-the-loop replays (support re-runs). Stripe's 24 hours
        is a pragmatic upper bound for interactive clients; batch
        reconcilers may want days.
  \item \textbf{A shared cache serves one user's account page to another.
        Root-cause it.}
        A personalized response reached a shared cache without
        \texttt{private}/\texttt{no-store}, or the key omitted the
        authorization dimension (missing \texttt{Vary: Authorization}).
        The fix is directive hygiene plus key review; the prevention is a
        policy lint: no 200 with \texttt{Set-Cookie} or per-user data may
        be \texttt{public}.
  \item \textbf{Optimistic versus pessimistic control for a hot inventory
        row --- decide.}
        Measure collision rate. Below a few percent, optimistic wins (no
        lock overhead, collisions pay the retry). At high collision rates
        the retry churn dominates: move to a per-aggregate queue or row
        lock, trading latency for throughput stability. There is no
        doctrinaire answer --- only a measured one.
  \item \textbf{Design the recovery path for a client that receives 409
        ``idempotency key in flight.''}
        Bounded poll with backoff (the server hints via
        \texttt{Retry-After}) until the twin completes and the replay
        becomes available; on expiry of the bound, surface an explicit
        ``outcome unknown, reconcile via GET'' state rather than guessing
        either success or failure.
\end{enumerate}

\subsection*{Staff (Q29--Q33)}
\begin{enumerate}[label=\textbf{Q\arabic*.}, leftmargin=1.6cm, start=29]
  \item \textbf{Design the end-to-end duplicate-defense architecture for a
        payment platform: client, edge, service, database.}
        Client mints one key per intent; edge terminates retry storms and
        never retries unsafe methods itself; service middleware maps the
        key protocol; the database enforces (account, key) uniqueness and
        commits effect+snapshot atomically; downstream consumers get
        effectively-once from the same keys propagated as message IDs.
        Each layer assumes the others fail.
  \item \textbf{Your CDN bill tripled after a ``cache everything''
        initiative. How do you find the waste?}
        Attribute misses by cause: unkeyed variance (accidental
        \texttt{Vary}-like inputs such as per-user query tokens),
        lifetime mis-tuning (one-second \texttt{max-age}), non-cacheable
        methods, and key explosion from ignored parameters. Then rebuild
        the policy matrix per surface --- immutable assets long-cached with
        content-hashed URLs, API responses short-cached with validators,
        personalized surfaces \texttt{private} or \texttt{no-store}.
  \item \textbf{When is serving stale data unacceptable, and what is the
        mechanism that enforces freshness under failure?}
        When a wrong answer is worse than no answer: balances, entitlement
        checks, safety interlocks. \texttt{must-revalidate} (and refusal
        of \texttt{stale-if-error}) converts origin failure into an
        explicit error rather than a confident lie; that is a product
        decision encoded in a header.
  \item \textbf{How do idempotency keys interact with event-driven
        consumers downstream of the API?}
        The key (or a derived operation ID) must travel with the emitted
        event as its dedup identity; consumers apply the same
        effectively-once pattern with their own dedup stores and windows.
        The API's replay store protects the synchronous boundary only ---
        the async boundary needs its own~\cite{kleppmann2017ddia}.
  \item \textbf{Governance: how do you make 40 teams adopt one idempotency
        and caching policy?}
        Codify it where it cannot be skipped: shared middleware for the key
        protocol (with conformance tests every service must pass), a
        gateway-level lint for \texttt{Cache-Control} legality, and golden
        dashboards for replay rate and hit ratio. Then make the paved road
        faster than the detour --- a platform library doing the right thing
        by default beats a policy document every
        time~\cite{gehl2021apidesign,nygard2018releaseit}.
\end{enumerate}

%% ----------------------------------------------------------------------------
\section{External Bibliography \& Further Reading}
\label{sec:ch10-bibliography}

\begin{itemize}
  \item \textbf{RFC 9110 --- HTTP Semantics}~\cite{rfc9110}. The normative
        source for method safety and idempotency, status codes 304/409/412/
        428, conditional request headers (\texttt{If-Match},
        \texttt{If-None-Match}, \texttt{If-Modified-Since}), ETag strength,
        and \texttt{Vary}.
  \item \textbf{RFC 9111 --- HTTP Caching.} The companion specification:
        \texttt{Cache-Control} directives, freshness lifetimes, heuristic
        freshness, \texttt{stale-while-revalidate} and
        \texttt{stale-if-error} (registered extension directives), and
        cache-invalidation-on-unsafe-methods. (Read alongside
        RFC~9110~\cite{rfc9110}.)
  \item \textbf{IETF Draft --- The Idempotency-Key HTTP Header Field.}
        The standardization effort behind the header: key generation
        guidance, fingerprinting, and the replay/mismatch/in-flight
        taxonomy implemented in Listing~10.1.
  \item \textbf{Stripe API Documentation --- Idempotent Requests.} The
        canonical production deployment: client-generated keys, 24-hour
        retention, response replay, and the parameter-mismatch error shape.
  \item \textbf{Kleppmann, \emph{Designing Data-Intensive
        Applications}}~\cite{kleppmann2017ddia}. Chapters on transactions
        and distributed systems: lost updates, serializability, fencing
        tokens, and the definitive treatment of why exactly-once is
        impossible and effectively-once is the real target.
  \item \textbf{Nygard, \emph{Release It!}}~\cite{nygard2018releaseit}.
        Stability patterns and war stories: how retries, timeouts, and
        caches behave under failure in real production systems.
  \item \textbf{Gehl, \emph{API Design Patterns}}~\cite{gehl2021apidesign}.
        Design-pattern treatment of idempotency, request identifiers, and
        standard methods as building blocks.
  \item \textbf{Google AIP-155 (Request identification) and AIP-194
        (Retries).} Google's API Improvement Proposals for request IDs and
        retry-safe design: which methods may be retried and how request
        identifiers make retries safe.
  \item \textbf{FastAPI and Pydantic documentation}~\cite{fastapidocs,pydanticdocs}.
        Middleware ordering, \texttt{TestClient} usage, and the validation
        behavior the chapter's listings rely on.
\end{itemize}

%% ----------------------------------------------------------------------------
\section{Capstone Projects}
\label{sec:ch10-capstones}

\subsection*{Capstone 10.1 --- Idempotency-Key Client SDK \hfill [junior]}
\textbf{Objective:} Build the client half of the protocol: a small Python
class wrapping \texttt{httpx} that mints one key per logical operation and
replays it across retries, speaking to Listing~10.1's server.

\textbf{Inputs/Datasets:} Listing~10.1 (server); the Northwind Robotics
transfer payload fixture from the test suite.

\textbf{Steps:} (1)~Implement \texttt{Operation} objects that mint a UUIDv4
key at construction. (2)~Implement \texttt{submit()} with bounded retry:
retry on connection error and on 409 (honoring \texttt{Retry-After}),
never on 422. (3)~Prove a retried submit yields exactly one ledger entry.
(4)~Prove a tampered payload under the same key surfaces the 422 to the
caller. (5)~Log every attempt with attempt number and outcome.

\textbf{Acceptance Criteria:}
\begin{itemize}
  \item[$\square$] One \texttt{Operation} retries up to 5 times and produces
        exactly one server-side mutation.
  \item[$\square$] 422 is never retried; 409 is retried with backoff.
  \item[$\square$] A second \texttt{Operation} mints a fresh key and
        executes a second transfer.
  \item[$\square$] All tests pass offline and deterministically.
\end{itemize}

\textbf{Stretch Goals:} Persist in-flight operations to a local file so a
crashed client resumes with the same key.

\subsection*{Capstone 10.2 --- Conditional-GET Sync Agent \hfill [junior]}
\textbf{Objective:} Write a catalog sync agent that polls
Listing~10.2's parts API on a fixed schedule and transfers bytes only when
representations actually changed.

\textbf{Inputs/Datasets:} Listing~10.2; a scripted mutation schedule (restock
events at chosen logical steps).

\textbf{Steps:} (1)~Keep per-SKU validator state (ETag + body) in a local
JSON file. (2)~On each poll, send \texttt{If-None-Match} when a validator
exists. (3)~On 304, keep the stored body and count a zero-byte sync; on 200,
replace state. (4)~Run 20 polls over the scripted schedule; report bytes
transferred versus a naive always-GET baseline. (5)~Prove a mid-run restock
is observed exactly once.

\textbf{Acceptance Criteria:}
\begin{itemize}
  \item[$\square$] 304 polls transfer zero representation bytes.
  \item[$\square$] The agent's local state always equals the server's
        current representation after each poll.
  \item[$\square$] The savings report is reproducible from the scripted
        schedule.
\end{itemize}

\textbf{Stretch Goals:} Add \texttt{Last-Modified}/\texttt{If-Modified-Since}
as a fallback when no ETag is present.

\subsection*{Capstone 10.3 --- Optimistic-Locking Order Service \hfill [mid]}
\textbf{Objective:} Extend Listing~10.3 into an order-reservation service in
which reserving stock and releasing stock are both conditional writes, with
a client that performs the read--decide--write--retry loop.

\textbf{Inputs/Datasets:} Listing~10.3; a seeded catalog of 10 SKUs; a
deterministic order script (seeded random choice of SKU and quantity).

\textbf{Steps:} (1)~Model reservations as quantity transitions guarded by
\texttt{If-Match}. (2)~Implement the client retry loop with bounded
attempts and deterministic (seeded) jitter. (3)~Drive 100 concurrent order
workers from the seeded script. (4)~Assert global conservation: final
quantity equals initial minus sum of accepted reservations. (5)~Emit a
contention report: attempts histogram, 412 rate per SKU.

\textbf{Acceptance Criteria:}
\begin{itemize}
  \item[$\square$] No lost updates under the 100-worker storm
        (conservation invariant holds exactly).
  \item[$\square$] Blind writes receive 428; stale writes receive 412.
  \item[$\square$] The contention report is identical across runs.
\end{itemize}

\textbf{Stretch Goals:} Add an adaptive escalation: SKUs whose 412 rate
exceeds a threshold switch to a serialized per-SKU queue; show the
throughput/latency trade-off in the report.

\subsection*{Capstone 10.4 --- Cache-Policy Workbench \hfill [senior]}
\textbf{Objective:} Turn Listing~10.5 into a policy evaluation tool: given a
traffic timeline and a library of candidate policies, rank them by origin
load, staleness window, and error resilience.

\textbf{Inputs/Datasets:} Listing~10.5; a synthetic one-hour request
timeline (request arrivals, origin mutation events, one scripted outage);
five candidate policies from \texttt{no-store} to
\texttt{max-age=300, stale-while-revalidate=60, stale-if-error=600}.

\textbf{Steps:} (1)~Load timelines and policies from a JSON fixture.
(2)~Simulate every (policy, timeline) pair. (3)~Produce a comparison table:
origin fetches, bytes saved, longest stale serve, requests failed during the
outage. (4)~Write a one-page recommendation memo choosing a policy per
surface (product page, pricing API, account page). (5)~Add a regression
test asserting the ranking is stable.

\textbf{Acceptance Criteria:}
\begin{itemize}
  \item[$\square$] The table is reproducible bit-for-bit from the fixtures.
  \item[$\square$] At least one policy is shown to serve zero failed
        requests during the scripted outage, and the memo explains why.
  \item[$\square$] The account page is assigned \texttt{no-store} or
        \texttt{private} with a one-sentence justification.
\end{itemize}

\textbf{Stretch Goals:} Model a shared cache with a \texttt{Vary} dimension
and demonstrate a misconfiguration that leaks across user segments.

\subsection*{Capstone 10.5 --- Correctness Audit of a Retry-Capable Gateway \hfill [staff]}
\textbf{Objective:} Design and execute a full duplicate-safety audit for a
fictitious API platform: enumerate every retry source (client SDK, proxy,
load balancer, queue redelivery), classify every mutation endpoint by
idempotency posture, and produce the remediation plan.

\textbf{Inputs/Datasets:} A synthetic platform registry you author (YAML: 20
endpoints with methods, retry sources, and current dedup posture);
Listings~10.1--10.5 as the reference implementation kit.

\textbf{Steps:} (1)~Build the retry-source inventory: for each endpoint,
list every layer that can duplicate a request. (2)~Classify endpoints:
spec-idempotent, key-protected, unprotected. (3)~Write a linter that flags
unprotected unsafe methods reachable by any retry source. (4)~For the top
three risk endpoints, specify the fix (key protocol adoption, PUT
redesign, or constraint-backed dedup) with migration steps. (5)~Define the
ongoing telemetry: replay rate, 409 rate, 422 rate, and the alert
thresholds that would have caught the war story of
Section~\ref{sec:ch10-warstory}.

\textbf{Acceptance Criteria:}
\begin{itemize}
  \item[$\square$] The registry linter flags exactly the seeded violations,
        with zero false positives on the protected endpoints.
  \item[$\square$] Every remediation references the mechanism that enforces
        it (unique constraint, fingerprint, fencing token) --- not just a
        header.
  \item[$\square$] The telemetry plan includes at least one alert tied to
        each failure mode in the war story.
  \item[$\square$] The plan states the blast radius of the idempotency-key
        store being unavailable and how the platform degrades.
\end{itemize}

\textbf{Stretch Goals:} Extend the audit across the async boundary: model
queue redelivery for each endpoint's emitted events and specify the
consumer-side dedup windows.

%% ----------------------------------------------------------------------------
\section{Python Dependencies Appendix}
\label{sec:ch10-deps}

All code runs offline after a one-time install. From PowerShell:

\begin{minted}{text}
python -m venv .venv
.\.venv\Scripts\Activate.ps1
pip install -r requirements.txt
\end{minted}

\begin{minted}{text}
fastapi>=0.110
pydantic>=2.7
httpx>=0.27
pytest>=8
\end{minted}

\texttt{sqlite3}, \texttt{hashlib}, \texttt{threading}, and
\texttt{dataclasses} are Python standard library and need no installation.

\subsection{fastapi ($\geq$ 0.110)}
\textbf{Description:} The ASGI web framework hosting every served API in
this chapter. \textbf{Why included:} routing, the middleware extension point
used by the idempotency layer (Listing~10.1), header access for conditional
requests (Listings~10.2, 10.3), and the application-factory pattern that
gives every test an isolated store~\cite{fastapidocs}.
\textbf{Alternatives:} Starlette directly (FastAPI's base); Flask plus
manual request validation. \textbf{Installation:}
\texttt{pip install "fastapi>=0.110"}.

\subsection{pydantic ($\geq$ 2.7)}
\textbf{Description:} Type-driven validation and serialization library.
\textbf{Why included:} defines the \texttt{TransferRequest} and
\texttt{InventoryUpdate} wire contracts; its declarative constraints
(\texttt{gt=0}, \texttt{ge=0}) are the boundary validation the correctness
arguments assume~\cite{pydanticdocs}. \textbf{Alternatives:}
\texttt{dataclasses} with hand-rolled checks; attrs; marshmallow.
\textbf{Installation:} \texttt{pip install "pydantic>=2.7"}.

\subsection{httpx ($\geq$ 0.27)}
\textbf{Description:} Synchronous and asynchronous HTTP client.
\textbf{Why included:} it is the transport underneath FastAPI's
\texttt{TestClient}; every offline test in this chapter runs through it, and
Capstone~10.1 builds the retrying client SDK on it. \textbf{Alternatives:}
requests (synchronous only, not ASGI-native); \texttt{urllib} (stdlib,
low-level). \textbf{Installation:} \texttt{pip install "httpx>=0.27"}.

\subsection{pytest ($\geq$ 8)}
\textbf{Description:} Test framework. \textbf{Why included:} runs all 33
tests in this chapter, including the barrier-synchronized concurrency
suites whose assertions prove exactly-once effect and lost-update
prevention. \textbf{Alternatives:} unittest (stdlib); nose2.
\textbf{Installation:} \texttt{pip install "pytest>=8"}.

%% ----------------------------------------------------------------------------
\section{Chapter Summary \& Key Takeaways}
\label{sec:ch10-summary}

\begin{itemize}
  \item \textbf{Idempotency is a fixed-point property} ---
        $f(f(x)) = f(x)$ on observable state. GET/PUT/DELETE have it by
        specification; POST gets it by design, via idempotency
        keys~\cite{rfc9110}.
  \item \textbf{The Idempotency-Key protocol is a small state machine:}
        client-minted key, server-side fingerprint (method + route +
        canonical body), atomic claim, response snapshot, replay on
        duplicate, 422 on mismatch, 409 while in flight, expiry matched to
        the client's retry horizon.
  \item \textbf{Exactly-once delivery is impossible} (Two Generals);
        at-least-once delivery plus idempotent processing is the
        effectively-once substitute, and every dedup mechanism has a window
        that must be sized and monitored~\cite{kleppmann2017ddia}.
  \item \textbf{Caching is a contract:} the origin assigns a freshness
        lifetime (\texttt{max-age}) and validators (strong ETag); the cache
        revalidates with \texttt{If-None-Match} and serves 304s;
        \texttt{no-cache} $\neq$ \texttt{no-store};
        \texttt{stale-while-revalidate} buys latency,
        \texttt{stale-if-error} buys availability, \texttt{must-revalidate}
        vetoes both.
  \item \textbf{Concurrency control is compare-and-set at the storage
        boundary.} Conditional writes with \texttt{If-Match} convert lost
        updates into 412s; 428 demands the precondition; 409 is for state
        conflicts. Distributed locks need fencing tokens or they are
        throughput hints, not correctness.
  \item \textbf{Retries without idempotency are unsafe by construction.}
        Every infrastructure layer that retries POST silently is a latent
        double-charge; the fix is architectural (keys, idempotent methods),
        not procedural (``be careful'').
  \item \textbf{Prove it, don't hope:} the chapter's tests collapse
        50-thread storms into single mutations and 20-writer races into
        exact arithmetic --- deterministic assertions on final state, never
        on timing.
\end{itemize}

%% ----------------------------------------------------------------------------
\section{Quick-Reference Card}
\label{sec:ch10-quickref}

\begin{table}[h]
  \centering
  \small
  \begin{tabular}{l c c l}
    \toprule
    \textbf{Method} & \textbf{Safe} & \textbf{Idempotent} & \textbf{Retry policy} \\
    \midrule
    GET / HEAD & yes & yes & always safe to retry \\
    OPTIONS / TRACE & yes & yes & always safe to retry \\
    PUT & no & yes & safe to retry \\
    DELETE & no & yes & safe to retry \\
    POST & no & no & retry only with an idempotency key \\
    PATCH & no & no & retry only if the delta is absolute or keyed \\
    \bottomrule
  \end{tabular}
  \caption{Method idempotency and retry policy~\cite{rfc9110}.}
\end{table}

\begin{table}[h]
  \centering
  \small
  \begin{tabular}{l l}
    \toprule
    \textbf{Directive} & \textbf{One-line meaning} \\
    \midrule
    \texttt{max-age=n} & serve fresh from cache for $n$ seconds \\
    \texttt{no-cache} & store OK, but revalidate before every reuse \\
    \texttt{no-store} & never store; always go to origin \\
    \texttt{private} / \texttt{public} & browser-only vs.\ shared-cacheable \\
    \texttt{stale-while-revalidate=n} & serve stale $\le n$s while refreshing \\
    \texttt{stale-if-error=n} & serve stale $\le n$s when origin errors \\
    \texttt{must-revalidate} & no stale serving once expired, ever \\
    \bottomrule
  \end{tabular}
  \caption{Cache-Control cheat-sheet.}
\end{table}

\begin{table}[h]
  \centering
  \small
  \begin{tabular}{l l}
    \toprule
    \textbf{Status} & \textbf{Select when} \\
    \midrule
    304 & validator still matches; no body retransmission needed \\
    409 & state conflict: business rule or in-flight idempotency twin \\
    412 & \texttt{If-Match}/precondition evaluated false (lost update refused) \\
    422 & idempotency key reused with a different payload fingerprint \\
    428 & server requires conditional writes; precondition header missing \\
    \bottomrule
  \end{tabular}
  \caption{Status-code selection for this chapter's machinery.}
\end{table}

%% ----------------------------------------------------------------------------
\section{Standardized Evidence Addendum}
\subsection{Benchmark Snapshot Template}
\begin{table}[h]
  \centering
  \small
  \begin{tabularx}{\textwidth}{l c c c c}
    \toprule
    \textbf{Config} & \textbf{p50 latency} & \textbf{p99 latency} & \textbf{Error rate} & \textbf{Throughput} \\
    \midrule
    Baseline & -- & -- & -- & 1.00x \\
    Candidate A & -- & -- & -- & -- \\
    Candidate B & -- & -- & -- & -- \\
    \bottomrule
  \end{tabularx}
  \caption{Minimum benchmark table to keep per chapter.}
\end{table}

\subsection{Request Trace Template}
\begin{itemize}
  \item \textbf{Request:} one representative API call for this chapter's topic.
  \item \textbf{Before:} behavior (headers, payload, latency, failure mode) with the naive approach.
  \item \textbf{After:} behavior with the technique in this chapter.
  \item \textbf{Result:} what changed in correctness, resilience, latency, and operability.
\end{itemize}

\subsection{Cost and Latency Budget Snapshot}
\begin{table}[h]
  \centering
  \small
  \begin{tabularx}{\textwidth}{l c c}
    \toprule
    \textbf{Stage} & \textbf{Latency budget} & \textbf{Cost driver} \\
    \midrule
    TLS / connection setup & -- & handshake round trips, keep-alive hit rate \\
    Request validation & -- & schema complexity, CPU per request \\
    Business logic and I/O & -- & downstream fan-out, query cost \\
    Serialization and egress & -- & payload size, compression ratio \\
    \bottomrule
  \end{tabularx}
  \caption{Operational budget template for this chapter's technique.}
\end{table}

\section{Configuration Preset Template}
\begin{minted}{text}
# api_policy.yaml
policy_version: "v1.0.0"
component: "<chapter_topic>"
timeout_ms: 0
retry_budget: 0
fallback_strategy: "fail_fast"
rollback_trigger: "error_budget_burn_gt_threshold"
\end{minted}

\section{CI Quality Gate Specification}
\begin{itemize}
  \item contract tests green against the pinned OpenAPI/AsyncAPI/Protobuf spec,
  \item no breaking schema changes without a version bump,
  \item error responses conform to RFC 9457 problem-details shape,
  \item benchmark deltas within approved regression bounds (latency and throughput),
  \item policy version and rollback plan present in release metadata.
\end{itemize}

\section{Incident Runbook Template}
\begin{enumerate}
  \item Detect regression from SLO burn-rate alerts or consumer reports.
  \item Scope impacted endpoints, versions, and consumer segments.
  \item Contain impact (rollback, feature flag, rate-limit tightening, or traffic shift).
  \item Diagnose with traces, structured logs, and contract diffs.
  \item Recover with validated fix and verified rollout procedure.
  \item Prevent recurrence via new regression tests and CI gates.
\end{enumerate}

\section{Cross-Chapter Decision Card}
\begin{table}[h]
  \centering
  \small
  \begin{tabularx}{\textwidth}{l X X}
    \toprule
    \textbf{Dimension} & \textbf{Choose this approach when} & \textbf{Prefer an alternative when} \\
    \midrule
    Use case & -- & -- \\
    Latency budget & -- & -- \\
    Operational cost & -- & -- \\
    Team maturity & -- & -- \\
    \bottomrule
  \end{tabularx}
  \caption{Decision card to complete per chapter.}
\end{table}

\section{ROI Model and Build-vs-Buy}
\begin{itemize}
  \item \textbf{Build when:} the capability is differentiating, compliance forbids vendors, or scale makes vendor pricing irrational.
  \item \textbf{Buy when:} the capability is commodity (auth, gateways, observability), time-to-market dominates, or staffing is thin.
  \item \textbf{Hidden costs to model:} on-call load, upgrade cadence, expertise ramp, data egress fees.
\end{itemize}

\section{Disaster Recovery Notes (RPO/RTO)}
\begin{itemize}
  \item Define recovery point objective (RPO): maximum acceptable data loss window.
  \item Define recovery time objective (RTO): maximum acceptable restoration time.
  \item Test restores regularly; an untested backup is not a backup.
\end{itemize}

